\chapter{Higher associativity and formality}
\label{chap:higher-associativity}

The relation $t^r=2uz$ holds in the cohomology of the central algebra.
We ask whether this differential graded algebra can be connected to
its cohomology by a finite zigzag of multiplicative quasi-isomorphisms.
A six-input operation will exclude such a zigzag for every $r\geq2$.

Let $k$ be a field of characteristic different from two, and let $r\geq1$ be an integer.
All tensor products in this chapter are over $k$, unless a coefficient
ring is displayed. Every element of a direct sum has finitely many
nonzero components.

Define the unital associative differential graded
algebra
\[
 C_r=k[t]\langle\tau\rangle[z]/(z^2),\qquad
 |t|=0,\quad |\tau|=-1,\quad |z|=2,
 \qquad D\tau=t^r-2\tau^2z,\quad Dt=Dz=0.
\]
The elements $t,z$ commute with every element. The differential has
degree one and obeys $D(ab)=(Da)b+(-1)^{|a|}aDb$.

Put $u=\tau^2$, so $|u|=-2$. Since $D\tau$ commutes with $\tau$,
\[
 Du=(D\tau)\tau-\tau D\tau=0.
\]
It follows that $D^2\tau=0$. The differential also preserves the
relations involving $t,z$, so these formulas define $C_r$.

A quasi-isomorphism is a degree-preserving algebra map that commutes
with differentials and induces an isomorphism on cohomology.
Formality over a coefficient ring means a finite zigzag of unital
quasi-isomorphisms to cohomology with zero differential. Every arrow
preserves the specified coefficient action. Formality over $k$ therefore
allows the images of $t,u,z$ to change.

\section{Normal forms and the first multiplication defect}
\label{sec:ha-normal-forms}

For $r=2$, the scalar $t^4$ has the primitive
\[
 w=\tau(t^2+2uz),\qquad Dw=t^4.
\]
Its square is also a boundary:
\[
 w^2=u(t^4+4t^2uz)
     =D\bigl(\tau u(t^2+6uz)\bigr).
\]
Thus $[w^2]=0$, so the square of this primitive gives no obstruction to formality.
The remaining question concerns compatible primitives for all
multiplication defects.

To choose primitives for all boundaries by a linear map, put
\[
 B=k[u,z]/(z^2),\qquad S=B[t],\qquad f=t^r-2uz.
\]
All elements of $S$ have even degree. Separating even and odd powers
of $\tau$ gives
\begin{equation}\label{eq:ha-two-summands}
 C_r=S\oplus\tau S,\qquad D(a+\tau b)=fb.
\end{equation}

The polynomial $f$ is monic, meaning that its leading coefficient is one.
Multiplication by $f$ is injective on $S$. Indeed, a nonzero polynomial
multiplied by the monic polynomial $f$ retains a nonzero leading
coefficient. This remains valid for coefficients in $B$, where $z^2=0$.
Consequently
\[
 H:=H^*(C_r)=k[u,t,z]/(t^r-2uz,z^2),
 \qquad H^{\mathrm{odd}}=0.
\]

Division by $f$ can be performed without dividing any coefficient in
$B$. If a polynomial has leading term $a_dt^d$ with $d\geq r$,
subtract $a_dt^{d-r}f$. The leading terms cancel, so the degree
strictly decreases. Repeat until the remaining polynomial has degree
less than $r$, or is zero. The degrees are nonnegative integers,
so this process terminates.

Adding the subtracted multiples gives a decomposition
\[
 a=fq+v,\qquad q,v\in B[t],\qquad
 v=0\ \text{or}\ \deg_t v<r.
\]
Suppose $a=fq'+v'$ is another such decomposition.
Then $f(q-q')=v'-v$. If $q-q'$ were nonzero, the left-hand side
would have degree at least $r$, by the leading-coefficient argument.
The right-hand side has degree less than $r$ or is zero.
Thus $q=q'$ and $v=v'$.

Adding two decompositions, or multiplying one by an element of $B$,
preserves the bound on the remainder's degree. Uniqueness therefore
makes both the quotient and remainder maps $B$-linear.
Since $f$ has cohomological degree zero, division also preserves that
grading. Only the leading coefficient one was used, so the same
construction applies to any monic polynomial over a commutative ring.

Every element of $H$ consequently has a unique representative of
$t$-degree less than $r$. Let $i:H\to S$ select it.
Let $p:C_r\to H$ take the class of the even component and kill the odd
component. These are degree-zero maps of complexes, and $pi=1_H$.

Define the degree-minus-one map
\[
 h(a)=\tau\frac{a-ip(a)}{f}\quad(a\in S),\qquad h(\tau a)=0.
\]
The numerator is divisible by $f$. On $S$, the equality $Dh=1-ip$
is its defining division identity. On $\tau S$, one has
$hD(\tau a)=h(fa)=\tau a$. Hence
\begin{equation}\label{eq:ha-contraction}
 Dh+hD=1-ip,\qquad h^2=hi=ph=0.
\end{equation}

Such maps are called a contraction onto cohomology. They choose
representatives and a primitive for the part removed by projection.
The linearity of division makes $i$ and $h$ $B$-linear, since
$\tau$ commutes with $B$. The quotient map $p$ is also $B$-linear.

The representative map $i$ generally fails to preserve products.
Its failure has the unique expression
\begin{equation}\label{eq:ha-F-definition}
 i(a)i(b)-i(ab)=f\,iF(a,b),
 \qquad F:H\otimes H\longrightarrow H.
\end{equation}
The quotient has $t$-degree less than $r$, because the product has
$t$-degree at most $2r-2$. Thus $F$ is defined without another choice.
It is $B$-bilinear and has degree zero.

For $r=2$, one has $H=B\oplus Bt$ and
\[
 F(bt,ct)=bc,\qquad F(B,H)=F(H,B)=0.
\]
In particular, $i(t)^2-i(t^2)=t^2-2uz=D\tau$.
Define the degree-minus-one map
\[
 f_2:H\otimes H\longrightarrow C_r,\qquad
 f_2(a,b)=-\tau iF(a,b).
\]
Then
\[
 Df_2(a,b)=i(ab)-i(a)i(b).
\]
For three inputs, the two product homotopies differ by
\[
 \begin{split}
 T(a,b,c)={}&i(a)f_2(b,c)-f_2(ab,c)\\
 &+f_2(a,bc)-f_2(a,b)i(c).
 \end{split}
\]

The displayed differential identity and the Leibniz rule give
\[
 \begin{split}
 DT={}&i(a)\bigl(i(bc)-i(b)i(c)\bigr)
       -\bigl(i(abc)-i(ab)i(c)\bigr)\\
 &+\bigl(i(abc)-i(a)i(bc)\bigr)
       -\bigl(i(ab)-i(a)i(b)\bigr)i(c)=0.
 \end{split}
\]
Equating the two homotopies requires a degree-minus-two map
$f_3:H^{\otimes3}\to C_r$ satisfying $Df_3=-T$.
Here $T$ lies in $\tau S$, where $D$ is injective, so $T=0$
and $f_3=0$ is a possible choice.
At four inputs, the comparison also contains
$f_2(a,b)f_2(c,d)$.
For $r=2$, its value on $(t,t,t,t)$ is $\tau^2=u$.

\section{Encoding multiplication and its successive homotopies}
\label{sec:ha-coalgebra}

An $n$-input map has arity $n$.
A tensor coalgebra records these compatibility equations together
with their signs. A degree shift makes each operation have degree one.

For a graded vector space $V$, define $sV$ by $(sV)^j=V^{j+1}$.
The map $s:V\to sV$ has degree $-1$, so $|sx|=|x|-1$.
The tensor convention for homogeneous maps is
\[
 (F\otimes G)(x\otimes y)=(-1)^{|G||x|}F(x)\otimes G(y).
\]
In particular,
\[
 s^{\otimes n}(x_1\otimes\cdots\otimes x_n)
 =(-1)^{\sum_{j=1}^n(n-j)|x_j|}
 sx_1\otimes\cdots\otimes sx_n.
\]

The tensor coalgebra is the vector space
\[
 T^c(sV)=\bigoplus_{n\geq1}(sV)^{\otimes n}
\]
with coproduct given by cutting a word into two nonempty consecutive
pieces. There is no length-zero summand or counit:
\[
 \Delta(v_1\otimes\cdots\otimes v_n)
 =\sum_{j=1}^{n-1}(v_1\otimes\cdots\otimes v_j)
        \otimes(v_{j+1}\otimes\cdots\otimes v_n).
\]
Cutting twice gives every decomposition into three pieces once in
either order. This is the coassociativity identity
$(\Delta\otimes1)\Delta=(1\otimes\Delta)\Delta$.
Let $\pi:T^c(sV)\to sV$ project to words of length one.
Composing a map with $\pi$ is called corestriction.

An element is primitive if its coproduct is zero.
Every word of length one is primitive.
For $n\geq2$, project the coproduct of a length-$n$ component
to $sV\otimes(sV)^{\otimes(n-1)}$.
This gives the original component, with its first factor separated.
Thus a primitive element has zero component in every length greater
than one, and
\begin{equation}\label{eq:ha-primitives}
 \ker\bigl(\Delta:T^c(sV)\to T^c(sV)\otimes T^c(sV)\bigr)=sV.
\end{equation}

A coderivation of degree $d$ is a map $b$ satisfying
$\Delta b=(b\otimes1+1\otimes b)\Delta$ with the tensor signs above.
It is determined by its components $b_n=\pi b|_{(sV)^{\otimes n}}$.
Their extension is
\begin{equation}\label{eq:ha-coderivation-extension}
 \begin{split}
 &b(v_1\otimes\cdots\otimes v_N)\\
 &\quad=\sum_{a+n+c=N}(-1)^{d(|v_1|+\cdots+|v_a|)}
 v_1\otimes\cdots\otimes v_a\otimes
 b_n(v_{a+1},\ldots,v_{a+n})\otimes
 v_{a+n+1}\otimes\cdots\otimes v_N .
 \end{split}
\end{equation}

Here $a,c\geq0$ and $n\geq1$.
An output cut lies before or after the inserted factor $b_n$.
It corresponds to an input cut before or after the interval replaced
by $b_n$. These terms are exactly
$(b\otimes1+1\otimes b)\Delta$, with the same prefix signs.
This proves the coderivation identity.

To prove uniqueness, suppose two coderivations have the same
corestrictions and agree on words shorter than $n$.
Their difference on a word of length $n$ has zero coproduct:
the coderivation identity expresses it using only differences on
shorter words. This difference is therefore primitive and has zero
corestriction. Equation~\eqref{eq:ha-primitives} makes it zero.
The same argument starts at length one, proving uniqueness by induction.

A degree-zero linear map $\Phi:T^c(sV)\to T^c(sW)$ is a
coalgebra map if it satisfies
\[
 \Delta_W\Phi=(\Phi\otimes\Phi)\Delta_V.
\]
Write $\pi_W:T^c(sW)\to sW$ for the length-one projection.
Its components are the corestrictions
\[
 \Phi_n=\pi_W\Phi|_{(sV)^{\otimes n}}:(sV)^{\otimes n}\to sW.
\]

Conversely, a family of such degree-zero components defines a
coalgebra map by summing over all consecutive partitions:
\begin{equation}\label{eq:ha-coalgebra-map}
 \Phi|_{(sV)^{\otimes N}}
 =\sum_{i_1+\cdots+i_q=N}\Phi_{i_1}\otimes\cdots\otimes\Phi_{i_q}.
\end{equation}
All indices are positive. A cut between two output factors is a cut
between two blocks, which proves compatibility with coproducts.
These formulas use finitely many intervals or partitions on each word.

Suppose coalgebra maps $\Phi,\Psi$ have the same components and
agree on words shorter than $n$. For a word $w$ of length $n$,
\[
 \Delta_W(\Phi-\Psi)(w)
 =\bigl((\Phi-\Psi)\otimes\Phi+
                \Psi\otimes(\Phi-\Psi)\bigr)\Delta_V(w)=0.
\]
Every cut of $w$ has two shorter pieces, so each term is zero.
The difference $(\Phi-\Psi)(w)$ is primitive with zero corestriction,
and is therefore zero by~\eqref{eq:ha-primitives}.
Starting at length one proves that the components determine the map.

\begin{definition}\label{def:ha-ainfinity}
An $A_\infty$-algebra on $V$ is a degree-one coderivation $b$ of
$T^c(sV)$ satisfying $b^2=0$. Its unshifted operations $m_n$ are
defined by
\[
 b_n s^{\otimes n}=s m_n,\qquad |m_n|=2-n.
\]

The structure is minimal if $m_1=0$. A morphism of these algebras
is a coalgebra map $\Phi$ satisfying $b^W\Phi=\Phi b^V$.
Its unshifted components $f_n$ satisfy
\[
 \Phi_n s^{\otimes n}=s f_n,\qquad |f_n|=1-n.
\]
It is a quasi-isomorphism if $f_1$ induces an isomorphism of the
cohomology groups for $m_1$.
No unit condition is imposed on the higher components.
\end{definition}

The degrees explain why multiplication requires a different test
from the linear obstructions of Chapter~\ref{chap:compatible-formal-lifts}.
For a mixed complex in that chapter, the maps
\[
 pB(hB)^{j-1}i:H_n\longrightarrow H_{n+2j-1},\qquad j\geq1,
\]
have odd chain degree $2j-1$.
If its homology is concentrated in even degrees, these maps vanish.
Here $m_{2j}:H^{\otimes2j}\to H$ has even cochain degree $2-2j$,
while $m_{2j+1}$ has odd degree.
Thus the evenness of $H$ forces the odd-arity operations to vanish
and permits the even-arity operations.

Corestricting $b^2=0$ and removing the shifts gives
\begin{equation}\label{eq:ha-stasheff}
 \sum_{a+j+c=n}(-1)^{a+jc}
 m_{a+1+c}(1^{\otimes a}\otimes m_j\otimes1^{\otimes c})=0.
\end{equation}

For an explicit sign calculation, put $d_i=|x_i|$ and
$S_n(d_1,\ldots,d_n)=\sum_i(n-i)d_i$.
The inner coderivation contributes $\sum_{i\leq a}(d_i-1)$.
The inserted unshifted output has degree
$d_{a+1}+\cdots+d_{a+j}+2-j$.
Adding the initial, inner, and outer shift exponents gives, modulo two,
\[
 \begin{split}
 &S_n(d)+\sum_{i\leq a}(d_i-1)
   +S_j(d_{a+1},\ldots,d_{a+j})\\
 &\quad+
 S_{a+1+c}\bigl(d_1,\ldots,d_a,
             \textstyle\sum_{i=a+1}^{a+j}d_i+2-j,
             d_{a+j+1},\ldots,d_n\bigr)\\
 &\hspace{2cm}\equiv a+jc+(2-j)\sum_{i\leq a}d_i .
 \end{split}
\]

The last term is the tensor-map sign for moving $m_j$ past the
first $a$ inputs. This proves the sign in~\eqref{eq:ha-stasheff}.
For $n=1,2$, this gives $m_1^2=0$ and the graded Leibniz rule.
For a minimal structure, $n=3$ gives associativity of $m_2$.

The same removal of shifts from $b^W\Phi=\Phi b^V$ gives
\begin{equation}\label{eq:ha-morphism}
 \begin{split}
 &\sum_{a+j+c=n}(-1)^{a+jc}
 f_{a+1+c}(1^{\otimes a}\otimes m_j^V\otimes1^{\otimes c})\\
 &\hspace{1cm}=
 \sum_{i_1+\cdots+i_q=n}
 (-1)^{\sum_{v=1}^{q-1}(q-v)(i_v-1)}
 m_q^W(f_{i_1}\otimes\cdots\otimes f_{i_q}).
 \end{split}
\end{equation}

For the right-hand side, let $d_v$ be the sum of the input degrees
in block $v$. The total shift exponent reduces to
\[
 \sum_{v<w}(1-i_w)d_v+
 \sum_{v=1}^{q-1}(q-v)(i_v-1).
\]
The first sum is the tensor-map sign for the maps $f_{i_w}$.
The second is therefore the remaining sign in the formula.

For a differential graded algebra, only $m_1=D$ and $m_2=\mu$ occur,
where $\mu$ is multiplication. Its shifted components are
\[
 b_1(sa)=sDa,\qquad b_2(sa,sb)=(-1)^{|a|}s(ab).
\]
Their square-zero equation states $D^2=0$, the Leibniz rule, and
associativity. Thus every algebra and map used in a formality zigzag
has this coalgebra description.

A differential graded algebra map gives a morphism whose only
nonzero component is $f_1$. Such a morphism is called strict.

\section{Transferring the product to cohomology}
\label{sec:ha-transfer}

The contraction~\eqref{eq:ha-contraction} chooses a primitive whenever
a product differs from its normal representative. Repeating this
operation constructs a minimal algebra on $H$.

\begin{proposition}\label{prop:ha-transfer}
Let $\mu$ denote multiplication in $C_r$. The recursive formulas
\begin{equation}\label{eq:ha-transfer-recursion}
 \begin{gathered}
 f_1=i,\qquad
 U_n=\sum_{a+b=n}(-1)^{a-1}\mu(f_a\otimes f_b),\\
 m_1=0,\qquad m_n=pU_n,\qquad f_n=-hU_n\quad(n\geq2)
 \end{gathered}
\end{equation}
define an $A_\infty$-algebra on $H$ and a quasi-isomorphism to $C_r$.
The operation $m_2$ is the cohomology product.
\end{proposition}

\begin{proof}
Suspend the contraction, writing $I=si s^{-1}$, $P=sp s^{-1}$,
$K=shs^{-1}$, and $D_1=sDs^{-1}$. Then
\[
 D_1K+KD_1=1-IP,\qquad PI=1,\qquad PK=KI=K^2=0.
\]
Let $D_1$ also denote its unary coderivation extension.
Let $\Lambda$ be the binary coderivation of multiplication in $C_r$.
The differential graded algebra identities give
\[
 D_1^2=\Lambda^2=0,\qquad D_1\Lambda+\Lambda D_1=0.
\]

Define, in increasing input length,
\[
 Q_n=\sum_{a+b=n}\Lambda_2(\Phi_a\otimes\Phi_b),\qquad
 \Phi_1=I,\quad\Phi_n=-KQ_n,\quad\beta_1=0,\quad\beta_n=PQ_n.
\]
Extend $\Phi$ as a coalgebra map and $\beta$ as a coderivation.
Write $Q=\pi\Lambda\Phi$. These definitions use only smaller
components at every stage.

Consider the defect $E=(D_1+\Lambda)\Phi-\Phi\beta$.
It satisfies
\[
 \Delta E=(E\otimes\Phi+\Phi\otimes E)\Delta.
\]

Suppose that $E$ vanishes on words shorter than $n$.
The displayed identity then makes $E$ of a word of length $n$
primitive. By~\eqref{eq:ha-primitives}, it has length one.
Thus $\pi\Lambda E=0$ on length $n$.

Substituting $E$, using $\Lambda^2=0$, and moving $\Lambda$
past $D_1$ with a minus sign gives
\[
 \pi\Lambda E=-D_1Q-Q\beta.
\]

The corestriction of the defect is therefore
\[
 \begin{split}
 \pi E
 &=-D_1KQ+Q-IPQ+KQ\beta\\
 &=K(D_1Q+Q\beta)=0
 \end{split}
\]
on length $n$. The defect already vanishes on length one.
Induction proves $(D_1+\Lambda)\Phi=\Phi\beta$.

The output component of greatest length of $\Phi$ is $I^{\otimes n}$.
It is injective because $P^{\otimes n}I^{\otimes n}=1$.
Taking the greatest nonzero input length proves that $\Phi$ is
injective on the direct sum. It follows that
\[
 \Phi\beta^2=(D_1+\Lambda)^2\Phi=0
 \quad\Longrightarrow\quad\beta^2=0.
\]
Unshifting a binary partition $(a,b)$ gives the sign $(-1)^{a-1}$
in~\eqref{eq:ha-morphism}, which proves
\eqref{eq:ha-transfer-recursion}. The linear component is $i$,
which induces the identity of $H$, so the morphism is a
quasi-isomorphism.
\end{proof}

Because $H$ has only even degrees, every odd-arity operation $m_n$
vanishes. The recursion gives more information about the even ones.
First,
\[
 f_2(a,b)=-\tau iF(a,b),\qquad f_3=0.
\]
Indeed, $U_2=i(a)i(b)$ and $-hU_2$ gives the first formula.
Every term of $U_3$ is odd and is killed by both $p$ and $h$.

For four inputs, the only surviving binary partition is $(2,2)$,
whose sign is negative. Therefore
\begin{equation}\label{eq:ha-m4}
 \begin{split}
 m_4(a,b,c,d)&=-uF(a,b)F(c,d),\\
 f_4(a,b,c,d)&=\tau u\,iF\bigl(F(a,b),F(c,d)\bigr).
 \end{split}
\end{equation}
Dividing $iF(a,b)iF(c,d)$ by $f$ gives the formula for $f_4$.

Induction in the recursion gives, for every $j\geq1$,
\[
 \begin{gathered}
 f_{2j+1}=0,\qquad
 f_{2j}(H^{\otimes2j})\subseteq\tau u^{j-1}S,\\
 m_{2j}(H^{\otimes2j})\subseteq u^{j-1}H.
 \end{gathered}
\]
In odd arity, each nonzero binary term has odd output and is killed
by $p$ and $h$.
In even arity greater than two, both surviving branches are
even-arity maps with odd output.
Their product contributes $\tau^2=u$, and the $B$-linearity of
$p,h$ preserves its coefficient.

For six inputs write $F_{ij}=F(a_i,a_j)$. The binary partitions
$(2,4)$ and $(4,2)$ each have negative sign, while $f_2$ also has
a negative sign. It follows that
\begin{equation}\label{eq:ha-m6}
 m_6(a_1,\ldots,a_6)
 =u^2\bigl(F_{12}F(F_{34},F_{56})
          +F(F_{12},F_{34})F_{56}\bigr).
\end{equation}

For $r=2$, the image of $F$ lies in $B$ and $F$ vanishes on a
$B$ input. Hence $f_4=0$ and all $f_n$ with $n\geq3$ vanish.
The only higher operation that remains is
\[
 m_4(t,t,t,t)=-u.
\]
Its nonzero value still does not decide formality.
A ternary change of coordinates on $H$ can remove it, and that
change modifies the sixth operation.

\section{A primitive for the fourth operation}
\label{sec:ha-hochschild}

Removing $m_4$ requires a ternary correction compatible with the
cohomology product. The Hochschild differential records this compatibility.

For integers $p\geq1$ and $q$, let
\[
 C^{p,q}(H,H)=\operatorname{Hom}_k^q(H^{\otimes p},H).
\]
The superscript $q$ means that the output degree is the sum of the
input degrees plus $q$. Since $H$ is even, these spaces are zero for
odd $q$. For a cochain $P\in C^{p,q}(H,H)$ define
\begin{equation}\label{eq:ha-hochschild-differential}
 \begin{split}
 (\delta P)(a_1,\ldots,a_{p+1})
 ={}&a_1P(a_2,\ldots,a_{p+1})\\
 &+\sum_{j=1}^{p}(-1)^j
 P(a_1,\ldots,a_ja_{j+1},\ldots,a_{p+1})\\
 &+(-1)^{p+1}P(a_1,\ldots,a_p)a_{p+1}.
 \end{split}
\end{equation}

All internal degrees here are even, so there are no additional element
signs. Applying $\delta$ twice merges each pair of intervals in both
orders with opposite signs. The overlapping merges agree by
associativity, and the end terms cancel in the same way.
Thus $\delta^2=0$. A cocycle satisfies $\delta P=0$, and a boundary
has the form $\delta Q$. Their quotient in bidegree $(p,q)$ is
the Hochschild cohomology group $HH^{p,q}(H,H)$.

Let $Q:M'\to M$ be an $A_\infty$-morphism between minimal
structures on the even space $H$ whose underlying coalgebra map
is invertible. Assume $q_1=1_H$ and $q_3$ is its first higher component.
The arity-two identity identifies their binary products.
The arity-four identity~\eqref{eq:ha-morphism} then gives
\[
 m'_4=m_4+\delta q_3.
\]
Thus removing $m_4$ requires $\delta q_3=-m_4$.

For cochains $P,Q$ of arities $p,\ell$, the cup product and insertion are
\[
 \begin{split}
 (P\smile Q)(a_1,\ldots,a_{p+\ell})
 &=P(a_1,\ldots,a_p)Q(a_{p+1},\ldots,a_{p+\ell}),\\
 (P\circ Q)(a_1,\ldots,a_{p+\ell-1})
 &=\sum_{j=0}^{p-1}(-1)^{j(\ell-1)}
 P(a_1,\ldots,a_j,Q(a_{j+1},\ldots,a_{j+\ell}),\ldots).
 \end{split}
\]
Expansion of~\eqref{eq:ha-hochschild-differential} gives
\[
 \delta(P\smile Q)=\delta P\smile Q+(-1)^pP\smile\delta Q.
\]
The middle exterior terms cancel. The remaining terms act within
the first block or within the second block, with the displayed signs.

Since $m_4=-uF\smile F$, a degree-minus-two unary cochain
$\theta$ with $\delta\theta=-2uF$ would give the primitive
$\tfrac12\theta\smile F$.
Differentiation in $z$ also detects the relation $z^2=0$, so both
defining relations enter the calculation.
Introduce central parameters $\alpha,\beta$ of degrees zero and
four, and the associative algebra
\[
 \mathcal H=k[u,t,z,\alpha,\beta]/
           (t^r-2uz-\alpha,z^2-\beta).
\]

Put $L=k[u,\alpha,\beta]$ and
$L_z=L[z]/(z^2-\beta)$. Division by the monic polynomial
$z^2-\beta$ makes $L_z$ free over $L$ on $1,z$.
Now write
\[
 \mathcal H=L_z[t]/(t^r-2uz-\alpha).
\]
The polynomial in $t$ is monic over the specified coefficient ring
$L_z$, so the same division construction makes $\mathcal H$ free
over $L_z$ on $1,t,\ldots,t^{r-1}$. Combining the two unique
expressions makes it free over $L$ on $t^iz^j$, where
$0\leq i<r$ and $j=0,1$.

Using this basis identifies its underlying module with $H[\alpha,\beta]$.
Let $\mu_{\alpha,\beta}$ be its multiplication in these coordinates.
The coefficients of $\alpha$ and $\beta$ in its expansion at zero are
\[
 F=\left.\partial_\alpha\mu_{\alpha,\beta}\right|_{0,0},
 \qquad G=\left.\partial_\beta\mu_{\alpha,\beta}\right|_{0,0}.
\]
These derivatives extract coefficients from a polynomial normal-form
product, so they are defined in every characteristic.

The first agrees with~\eqref{eq:ha-F-definition}. Their internal
degrees are zero and $-4$. The coefficients of $\alpha$ and $\beta$
in associativity give respectively
\begin{equation}\label{eq:ha-FG-cocycles}
 \delta F=0,\qquad\delta G=0.
\end{equation}
For example, the first coefficient equates
$F(ab,c)+F(a,b)c$ with $F(a,bc)+aF(b,c)$,
which is exactly $\delta F=0$.

For $a\geq0$, $0\leq i<r$, and $j\in\{0,1\}$, define
coefficient derivatives on the normal basis by
\[
 \theta_u(u^at^iz^j)=a u^{a-1}t^iz^j,
 \qquad
 \theta_z(u^at^iz^j)=j u^at^iz^{j-1}.
\]
The first expression is zero when $a=0$, and the second is zero
when $j=0$.
The factors $a,j$ denote the images of these integers in $k$.
Their degrees are $2$ and $-2$. These formulas differentiate the chosen
representatives, so the maps need not preserve products in $H$.

Their defects can be computed without differentiating a quotient.
For $x=u^at^iz^j$, $y=u^bt^\ell z^m$, put
$A=a+b$, $N=i+\ell$, and $J=j+m$. Then
\[
 F(x,y)=
 \begin{cases}
 u^At^{N-r}z^J,&N\geq r,\ J\leq1,\\
 0,&\text{otherwise},
 \end{cases}
\]
and
\[
 G(x,y)=
 \begin{cases}
 u^At^N,&N<r,\ J=2,\\
 2u^{A+1}t^{N-r}z^{J-1},&N\geq r,\ J\geq1,\\
 0,&\text{otherwise}.
 \end{cases}
\]
The inequality $N\leq2r-2$ allows only one replacement of $t^r$.
The resulting $z$-power needs at most one replacement of $z^2$.
These two replacements prove the formulas for $F,G$.

They also give
\begin{equation}\label{eq:ha-derivative-defects}
 \delta\theta_u=-2zF,
 \qquad\delta\theta_z=-2uF+2zG.
\end{equation}
For the first identity, differentiating the factor $2u$ introduced
by $t^r=2uz$ gives the extra $2zF$ in the derivative of a product.
The differential $\delta$ subtracts that derivative, giving $-2zF$.
For the second identity, differentiation of the same replacement
gives $-2uF$. Differentiating the two factors of a product killed
by $z^2=0$ contributes $2zG$.

The overlap of the replacements can be checked in the following
complete list of nonzero defects:
\[
 \begin{array}{c|c|c|c}
 \text{condition}&J&(\delta\theta_u)(x,y)&(\delta\theta_z)(x,y)\\\hline
 N<r&2&0&2u^At^Nz\\
 N\geq r&0&-2u^At^{N-r}z&-2u^{A+1}t^{N-r}\\
 N\geq r&1&0&2u^{A+1}t^{N-r}z
 \end{array}.
\]
All other cases give zero in both columns.
Substitution of the formulas for $F,G$ gives exactly these entries,
proving~\eqref{eq:ha-derivative-defects} on the full normal basis.

For these two cocycles, insertion gives
\[
 (F\circ G)(a,b,c)=F(G(a,b),c)-F(a,G(b,c)),
\]
and
\begin{equation}\label{eq:ha-insertion-FG}
 \delta(F\circ G)=G\smile F-F\smile G.
\end{equation}
To verify the sign, expand the five terms of the Hochschild
differential of $F\circ G$:
\[
 \begin{split}
 \delta(F\circ G)(a,b,c,d)
 ={}&aF(G(b,c),d)-aF(b,G(c,d))\\
 &-F(G(ab,c),d)+F(ab,G(c,d))\\
 &+F(G(a,bc),d)-F(a,G(bc,d))\\
 &-F(G(a,b),cd)+F(a,G(b,cd))\\
 &+F(G(a,b),c)d-F(a,G(b,c))d.
 \end{split}
\]

Apply $\delta F=0$ to the triples $(a,b,G(c,d))$ and
$(G(a,b),c,d)$. The resulting exterior products are
$G(a,b)F(c,d)-F(a,b)G(c,d)$.
The remaining terms with outer $F(-,d)$ combine using
$G(ab,c)-G(a,bc)=aG(b,c)-G(a,b)c$.
Those with outer $F(a,-)$ combine using the same identity on
$(b,c,d)$. The remaining four terms are the equation
$\delta F(a,G(b,c),d)=0$.
This proves~\eqref{eq:ha-insertion-FG}.

Define the arity-three, degree-minus-two cochain
\begin{equation}\label{eq:ha-g-definition}
 g=\tfrac12(\theta_z\smile F+\theta_u\smile G)-z(F\circ G).
\end{equation}
Combining the cup differential with
\eqref{eq:ha-derivative-defects} and~\eqref{eq:ha-insertion-FG} gives
\[
 \begin{split}
 \delta g
 &=-uF\smile F+z(G\smile F-F\smile G)
                  -z(G\smile F-F\smile G)\\
 &=-uF\smile F=m_4.
 \end{split}
\]
Thus the fourth operation is a Hochschild boundary for every $r$.
The primitive is $k$-linear, as permitted by formality over $k$.

The formulas for $F,G,\theta_u,\theta_z$ have integer coefficients
on the normal basis, and the cup and insertion identities are signed
multilinear expansions.
The cochain $g$ uses only $2^{-1}$.
Thus $\delta g=m_4$ also holds when the characteristic divides $r$.

\section{The sixth operation after changing coordinates}
\label{sec:ha-coordinate-change}

Let $M$ denote the transferred structure on $H$. Define a coalgebra
map $Q:T^c(sH)\to T^c(sH)$ by its unshifted components
\[
 q_1=1_H,\qquad q_3=-g,\qquad q_n=0\quad(n\ne1,3).
\]
The difference $Q-1$ lowers tensor length by at least two.
Hence
\[
 Q^{-1}=\sum_{j\geq0}(-1)^j(Q-1)^j
\]
is a finite sum on every input. Its inverse also preserves the
coproduct, by applying $Q^{-1}\otimes Q^{-1}$ to the coalgebra-map
identity.

Conjugation $b'=Q^{-1}bQ$ therefore gives a square-zero
coderivation, hence a new minimal structure $M'$ on $H$.
The map $Q:M'\to M$ is an isomorphism of these structures.

The augmentation is the graded algebra map
\[
 \epsilon:H\longrightarrow k,\qquad
 \epsilon(t)=\epsilon(u)=\epsilon(z)=0.
\]
The defining relations map to zero, so the map is well-defined.

The full space $H$ is needed for $Q$.
Indeed, $u,z\in\ker\epsilon$, whereas $q_3(u,z,z)=-1/2$.
Thus $Q$ does not restrict to $T^c(s\ker\epsilon)$.

Equation~\eqref{eq:ha-morphism} in arity four gives
\[
 m'_4=m_4+\delta q_3=m_4-\delta g=0.
\]
The arity-two morphism identity gives $m'_2=m_2$.
Degree parity gives $m'_3=m'_5=0$.

In arity six, the source-side terms containing $q_3$ also contain
$m'_4$, and therefore vanish. The four target partitions with one
block of size three and three blocks of size one all have positive
sign. In fact, their sign exponents are $2(3-j)$ for
$j=0,1,2,3$. The partition $(3,3)$ has exponent $2$, also even.
It follows that
\begin{equation}\label{eq:ha-six-coordinate}
 m'_6=m_6-\sum_{j=0}^3
 m_4(1^{\otimes j}\otimes g\otimes1^{\otimes(3-j)})+g\smile g.
\end{equation}

Equation~\eqref{eq:ha-m6} gives $\epsilon m_6=0$.
Every term with outer operation $m_4$ has a factor $u$, so
\begin{equation}\label{eq:ha-six-augmentation}
 \epsilon m'_6=(\epsilon g)\smile(\epsilon g).
\end{equation}

The arity-seven identity~\eqref{eq:ha-stasheff} reduces to
$-\delta m'_6=0$. Thus $m'_6$ defines a class in $HH^{6,-4}(H,H)$.

\section{A bar cycle detects the sixth operation}
\label{sec:ha-bar-pairing}

To annihilate every Hochschild boundary at once, we need a chain
whose adjacent products cancel. Put $I=\ker\epsilon$.
The reduced bar complex has underlying graded vector space
\[
 \mathcal B(H)=\bigoplus_{n\geq0}(sI)^{\otimes n}.
\]

The length-zero term is $k$. Write
$[a_1|\cdots|a_n]$ for $sa_1\otimes\cdots\otimes sa_n$.
Its degree is $\sum_i|a_i|-n$.
Every element of $I$ has even degree, so every shifted letter is odd.
The differential is consequently
\begin{equation}\label{eq:ha-bar-differential}
 b[a_1|\cdots|a_n]
 =\sum_{j=1}^{n-1}(-1)^{j-1}
 [a_1|\cdots|a_ja_{j+1}|\cdots|a_n].
\end{equation}

The product remains in $I$ because $\epsilon$ is multiplicative.
Two disjoint merges occur with opposite signs in $b^2$.
Two overlapping merges agree by associativity and also have opposite
signs. Hence $b^2=0$.

For two words, a shuffle is an interleaving that preserves the order
within each original word. The shuffle product is the sum of these
interleavings with the sign of the permutation of their shifted letters.
Here that sign is $(-1)^\nu$, where $\nu$ counts pairs moved past
each other. Denote the product by $\amalg$.

Since $H$ is commutative and even,
\[
 b(x\amalg y)=bx\amalg y+(-1)^{|x|}x\amalg by.
\]
Indeed, a merge within one source word gives the corresponding
term on the right. A merge between different source words appears
twice, with reversed adjacent letters and opposite signs.
Commutativity identifies their products, so these terms cancel.

Assume $r\geq2$, and put $A=t^{r-1}$. All letters in
\[
 \rho=[A|t]-[u|z]-[z|u],\qquad \sigma=[z|z]
\]
belong to $I$. Their differentials are
\[
 b\rho=t^r-uz-zu=0,\qquad b\sigma=z^2=0.
\]

Thus
\begin{equation}\label{eq:ha-detecting-chain}
 c=\rho\amalg\rho\amalg\sigma
\end{equation}
is a cycle of length six. Its internal degree, the sum of the
unshifted input degrees, is four. Its total bar degree is $4-6=-2$.
The degree-minus-four cochain $m'_6$ evaluates it in $H^0$.

Write $\gamma=\epsilon g$. The outer factor $z$ in
\eqref{eq:ha-g-definition} disappears under augmentation.
On normal monomials, $\epsilon\theta_z$ is nonzero only at $z$,
where it is one. Similarly, $\epsilon\theta_u$ is nonzero only at $u$.

The constant coefficient of $F(t^i,t^j)$ is one exactly when
$i+j=r$. The constant coefficient of $G(x,y)$ is one only at
$(x,y)=(z,z)$. Therefore, among the letters used by $c$, the only
possibly nonzero triples for $r>2$ are
\[
 Z_1=[z|A|t],\qquad Z_2=[z|t|A],\qquad U=[u|z|z].
\]
Each has $\gamma$-value $1/2$.

The signed coefficients in $c$ on concatenations of these triples
are
\[
 \begin{array}{c|rrr}
       &Z_1&Z_2&U\\\hline
 Z_1&2&0&-2\\
 Z_2&0&0&0\\
 U&-2&0&0
 \end{array}.
\]

Here is an explicit enumeration of the nonzero entries.
Abbreviate the three source words of $\rho$ by
$\alpha=[A|t]$, $\eta=[u|z]$, and $\xi=[z|u]$.
Their coefficients are $1,-1,-1$.
For each ordered pair of source words, interleave their letters with
$[z|z]$, preserving the order in all three words.
Before multiplying by the source coefficients, the nonzero counts are
\[
 \begin{array}{c|rrr}
 \text{ordered sources}&Z_1Z_1&Z_1U&UZ_1\\\hline
 (\alpha,\alpha)&2&0&0\\
 (\alpha,\eta),(\eta,\alpha)&0&0&1\ \text{each}\\
 (\alpha,\xi),(\xi,\alpha)&0&1\ \text{each}&0
 \end{array}.
\]

Each count is the sum of $(-1)^\nu$ over placements of the two
letters of each source. For $Z_1Z_1$, the two choices assigning the
two $\alpha$ words both have even $\nu$.
For $Z_1U$, the choices involving $\eta$ cancel, while each choice
involving $\xi$ has signed sum one.
For $UZ_1$, each choice involving $\eta$ has two positive placements
and one negative placement. Those involving $\xi$ have no admissible placement.

For $UU$, the pair $(\eta,\eta)$ has six placements of each sign.
Each mixed pair $(\eta,\xi)$ or $(\xi,\eta)$ has three of each sign.
The pair $(\xi,\xi)$ has no placement, because its words start with $z$.

A word containing $Z_2$ places $t$ before all available $A$ letters,
or ends in $A$, contrary to the order in $\alpha$.
These possibilities exhaust the nine ordered pairs of source words.

For $r=2$, put $Z=[z|t|t]$.
In $ZZ$, the source $[z|z]$ occupies positions one and four.
Choosing two of the remaining four positions for the first ordered
$[t|t]$ gives six placements.
For the subsets
\[
 (2,3),\ (2,5),\ (2,6),\ (3,5),\ (3,6),\ (5,6),
\]
the signs are respectively
\[
 +,\ -,\ +,\ +,\ -,\ +.
\]
Thus the coefficient of $ZZ$ is two.

The coefficients of $ZU,UZ,UU$ retain the preceding counts:
each word with $U$ contains at most one source $[A|t]$, whose
two positions still have their prescribed order after $A=t$.
The matrix on $Z,U$ is therefore
\[
 \begin{pmatrix}2&-2\\-2&0\end{pmatrix}.
\]
Thus for every $r\geq2$,
\begin{equation}\label{eq:ha-nonzero-pairing}
 (\epsilon m'_6)(c)=(\gamma\smile\gamma)(c)
 =\tfrac14(2-2-2)=-\tfrac12.
\end{equation}
This value is nonzero in $k$ because $2\ne0$.

If $m'_6=\delta K$ for $K\in C^{5,-4}(H,H)$, augmentation kills
the two exterior terms of~\eqref{eq:ha-hochschild-differential}
on inputs from $I$. The interior terms have the opposite signs
to~\eqref{eq:ha-bar-differential}. Hence
\[
 (\epsilon\delta K)(c)=-(\epsilon K)(bc)=0.
\]
This contradicts~\eqref{eq:ha-nonzero-pairing}. We have proved
\begin{equation}\label{eq:ha-nonzero-class}
 [m'_6]\ne0\quad\text{in }HH^{6,-4}(H,H).
\end{equation}

The pairing uses inputs from $I$.
Hochschild boundaries and coordinate changes are computed on the
full algebra $H$.

\section{Transport through arbitrary quasi-isomorphisms}
\label{sec:ha-transport}

To exclude a formality zigzag, we must transport the minimal model
across quasi-isomorphisms in either direction.
Only arities through six enter the obstruction.
The comparison will also be needed over coefficient rings, where an
acyclic kernel need not have a contracting homotopy.
Projectivity of the source will supply the individual lifts that the
construction requires.

For the lifting and transport constructions in this section, let $S$
be any commutative unital ring.
An $S$-algebra here is associative and unital, and the specified map
from $S$ has central degree-zero image killed by the differential.
All maps preserve this action, and all tensor products in these
constructions are over $S$.
The previous coalgebra definitions apply to graded $S$-modules with
the same tensor signs and algebraic direct sums.

The proofs of the component extension formulas use only consecutive
partitions and signed sums, so they still prove the coalgebra and
coderivation identities over $S$.
The primitive-element argument also remains valid: the cut after
the first letter is the reassociation isomorphism
$(sV)^{\otimes n}\cong sV\otimes(sV)^{\otimes(n-1)}$ on each
length $n\geq2$.
Thus the same primitive induction proves component uniqueness.
An $A_\infty$-structure over $S$ means a square-zero degree-one
coderivation of this tensor coalgebra, with the component conventions
in Definition~\ref{def:ha-ainfinity}.
Minimality means $b_1=0$, and a morphism is an $S$-linear coalgebra
map satisfying $b^W\Phi=\Phi b^V$.

For a positive integer $N$, write
\[
 T^c_{\leq N}(sV)=\bigoplus_{1\leq n\leq N}(sV)^{\otimes n}.
\]
Coderivations and coalgebra maps restrict to this space because their
formulas do not increase tensor length. An equation through arity $N$
means the corresponding equation on this truncated coalgebra.

An $S$-module $M$ is projective if every $S$-linear map $M\to Y$
lifts through every surjective $S$-linear map $X\to Y$.
A free module has this property: choose a preimage in $X$ of the
image of each basis vector, then extend linearly.
A direct summand of a free module is projective, by extending the
map to the free module using the projection and restricting its lift.
Conversely, a projective module is a direct summand of a free module.
Indeed, the free module on its underlying set surjects onto it, and
lifting its identity gives a section of that surjection.

Finite tensor products of projective modules are projective.
For splittings $M_i\xrightarrow{j_i}F_i\xrightarrow{p_i}M_i$
with $p_ij_i=1$, tensoring gives
$(\bigotimes_i p_i)(\bigotimes_i j_i)=1$.
The tensor product of the free modules $F_i$ is free on tuples of
basis vectors, so $\bigotimes_i M_i$ is its direct summand.
A direct sum of projective modules is also projective, since a map
from it lifts by lifting its restriction to each summand.

Say that a graded module is degreewise projective if each of its
homogeneous components is projective.
For such a $V$, put $W_n=(sV)^{\otimes_S n}$.
Each component
\[
 W_n^j=\bigoplus_{j_1+\cdots+j_n=j}
 (sV)^{j_1}\otimes_S\cdots\otimes_S(sV)^{j_n}
\]
is projective by the preceding two constructions.
These sums may have infinitely many summands, but each element has
finite support.
Consequently a degree-zero map from $W_n$ lifts through a degreewise
surjective graded map by lifting separately in each degree.
Over a field every vector space is free, so these conclusions include
the earlier field setting.

\begin{lemma}\label{lem:ha-surjective-lifting}
Let $q:P\to A$ be a degreewise surjective quasi-isomorphism of
differential graded $S$-algebras.
Let $(V,b)$ be a minimal $A_\infty$-algebra over $S$ whose
underlying graded module is degreewise projective.
Every coalgebra chain map
\[
 F:T^c_{\leq N}(sV)\longrightarrow T^c_{\leq N}(sA)
\]
lifts through $q$ to a coalgebra chain map
$\Phi:T^c_{\leq N}(sV)\to T^c_{\leq N}(sP)$.
If $F_1$ is a quasi-isomorphism, so is $\Phi_1$.
\end{lemma}

\begin{proof}
The complex $L=\ker q$ is acyclic. Take a cycle
$x\in L$. Injectivity of $H(q)$ gives $x=dy$ in $P$.

The element $qy$ is a cycle. Surjectivity of $H(q)$ gives a cycle
$z\in P$ and an element $a\in A$ such that $qy-qz=da$.
Choose $w\in P$ with $qw=a$. Then $y-z-dw\in L$ and
$d(y-z-dw)=x$, proving the assertion.

Suppose the lift $\Phi$ is constructed through length $n-1$.
The source $W_n=(sV)^{\otimes_S n}$ is degreewise projective,
so choose a degree-zero lift $\widetilde\Phi_n:W_n\to sP$ of $F_n$.
Use it as the new component of $\Phi$, and extend the components
by consecutive partitions. Let
\[
 E=D_P\Phi-\Phi b,
\]
where $D_P$ is the coderivation of the target differential and product.
The defect satisfies
\[
 \Delta E=(E\otimes\Phi+\Phi\otimes E)\Delta,
 \qquad D_PE+Eb=0.
\]

The first identity follows from the coalgebra identities for
$D_P,\Phi,b$. The second follows by expanding and using
$D_P^2=b^2=0$.

Since the defect vanishes below length $n$, its value in length $n$
is primitive. It lies in $sP$ by~\eqref{eq:ha-primitives}.
Minimality means that
$b$ lowers length, so $Eb=0$ on this input length.
It follows that $D_{P,1}E_n=0$.

The equality $q\Phi=F$ through length $n$ also gives $qE_n=0$.
Write $Z^j(sL)=\ker(D_{P,1}:(sL)^j\to(sL)^{j+1})$.
In each degree the defect is a map
\[
 E_n^j:W_n^j\longrightarrow Z^{j+1}(sL).
\]
Acyclicity of $L$ makes
\[
 D_{P,1}:(sL)^j\longrightarrow Z^{j+1}(sL)
\]
surjective.
Projectivity of $W_n^j$ supplies a map $T_n^j:W_n^j\to(sL)^j$
with $D_{P,1}T_n^j=E_n^j$.

Assembling the degree components gives $T_n:W_n\to sL$ of degree zero.
Replace the new component by
\[
 \Phi_n=\widetilde\Phi_n-T_n.
\]
Its image under $q$ is unchanged.
All contributions involving products or $b$ in input length $n$
use lower components of $\Phi$.
Consequently the replacement changes the defect only by
$-D_{P,1}T_n=-E_n$ and removes it in that length.

The correction lifts this particular map from $W_n$ into cycles.
It does not require a contracting homotopy on $L$.
The same argument starts at length one and proves the lift by induction.
Finally, $H(q)H(\Phi_1)=H(F_1)$ proves the assertion about the
linear component.
\end{proof}

A backward arrow need not be surjective. Its replacement will use
an algebra with two evaluations and a contraction onto the first.
The first evaluation imposes compatibility with the given arrow.
The second evaluation will provide the required surjection.

Define a graded free $S$-module
\[
 J^0=Se_0\oplus Se_1,\qquad J^1=S\eta,\qquad
 J^n=0\quad(n\ne0,1).
\]
Give it the bilinear product
\[
 \begin{split}
 &(ae_0+be_1+c\eta)(a'e_0+b'e_1+c'\eta)\\
 &\qquad=aa'e_0+bb'e_1+(ac'+cb')\eta,
 \qquad a,b,c,a',b',c'\in S.
 \end{split}
\]
For three factors, either parenthesization gives coefficients
$aa'a''$, $bb'b''$ on $e_0,e_1$ and coefficient
$aa'c''+ac'b''+cb'b''$ on $\eta$.
Thus the product is associative, with unit $e_0+e_1$.
It respects degrees because products of two degree-one elements are zero.

The product gives
\[
 e_0\eta=\eta,\qquad \eta e_1=\eta,\qquad
 \eta e_0=e_1\eta=\eta^2=0.
\]
Define a degree-one differential by
\[
 d_Je_0=-\eta,\qquad d_Je_1=\eta,\qquad d_J\eta=0.
\]
On a homogeneous element $x$, these formulas are equivalent to
\[
 d_Jx=\eta x-(-1)^{|x|}x\eta.
\]

For homogeneous $x,y$, expansion gives
\[
 \begin{split}
 (d_Jx)y+(-1)^{|x|}x\,d_Jy
 &=
 \eta xy-(-1)^{|x|}x\eta y
 +(-1)^{|x|}x\eta y
 -(-1)^{|x|+|y|}xy\eta\\
 &=d_J(xy).
 \end{split}
\]
The same expression for the differential gives
$d_J^2x=\eta^2x-x\eta^2=0$.
Also $d_J(e_0+e_1)=0$, so $J$ is a unital differential graded algebra.

Define evaluations $\epsilon_0,\epsilon_1:J\to S$ and a constant inclusion by
\[
 \begin{split}
 \epsilon_0(ae_0+be_1+c\eta)&=a,\\
 \epsilon_1(ae_0+be_1+c\eta)&=b,\\
 \iota:S\longrightarrow J,\qquad \iota(c)&=c(e_0+e_1).
 \end{split}
\]
The product formula makes each evaluation multiplicative.
Both evaluations kill $d_J(J)$ and send the unit to one.
The map $\iota$ preserves the product and has closed image.
Thus these are unital differential graded algebra maps, and
$\epsilon_0\iota=\epsilon_1\iota=1_S$.

Define the degree-minus-one map $K:J\to J$ by
\[
 K(\eta)=e_1,\qquad K(e_0)=K(e_1)=0.
\]
The maps $d_JK+Kd_J$ and $1-\iota\epsilon_0$ both take the values
$-e_1,e_1,\eta$ on $e_0,e_1,\eta$, respectively.
Consequently
\[
 d_JK+Kd_J=1-\iota\epsilon_0,\qquad
 \epsilon_0K=K\iota=K^2=0.
\]
Together with $\epsilon_0\iota=1_S$, these equations give the required
contraction onto the first endpoint.
The displayed basis formulas have integer coefficients, and no division occurs.

For a differential graded $S$-algebra $A$, give $A\otimes_S J$ the
differential and product
\[
 \begin{split}
 d(a\otimes x)&=da\otimes x+(-1)^{|a|}a\otimes d_Jx,\\
 (a\otimes x)(b\otimes y)&=(-1)^{|x||b|}ab\otimes xy
 \end{split}
\]
on homogeneous factors.
For a product of three tensors, the two parenthesizations give the same
tensor $abc\otimes xyz$.
Their sign exponents agree because
\[
 |x||b|+(|x|+|y|)|c|
 =|y||c|+|x|(|b|+|c|).
\]
Thus multiplication is associative, with unit
$1_A\otimes(e_0+e_1)$.

To check the differential, put $p=|a|$, $q=|b|$, and $\ell=|x|$.
Applying it twice gives
\[
 \begin{split}
 d^2(a\otimes x)
 ={}&d_A^2a\otimes x
 +\bigl((-1)^{p+1}+(-1)^p\bigr)da\otimes d_Jx
 +a\otimes d_J^2x\\
 ={}&0.
 \end{split}
\]
For the Leibniz rule, both
$d((a\otimes x)(b\otimes y))$ and
$d(a\otimes x)(b\otimes y)+(-1)^{p+\ell}(a\otimes x)d(b\otimes y)$
expand to
\[
 \begin{split}
 (-1)^{\ell q}\bigl(&
 (da)b\otimes xy+(-1)^p a(db)\otimes xy\\
 &+(-1)^{p+q}ab\otimes(d_Jx)y
 +(-1)^{p+q+\ell}ab\otimes x\,d_Jy\bigr).
 \end{split}
\]
The last two signs use the Leibniz rule in $J$.
These calculations prove that $A\otimes J$ is a differential graded algebra.

The endpoint maps and constant map extend to
\[
 \operatorname{ev}_i=1\otimes\epsilon_i:A\otimes J\longrightarrow A,
 \qquad
 \operatorname{const}(a)=a\otimes(e_0+e_1),\qquad i=0,1.
\]
They preserve units and commute with the differentials.
Each $\epsilon_i$ vanishes in degree one, so every surviving
tensor-product sign in the endpoint multiplication equation is positive.
Together with multiplicativity of $\epsilon_i$, this proves that
$\operatorname{ev}_i$ is multiplicative.
The degree-zero unit of $J$ makes $\operatorname{const}$ multiplicative.

The contraction also extends, with the tensor sign
\[
 K_A(a\otimes x)=(-1)^{|a|}a\otimes Kx.
\]
For homogeneous $a$, the two composites are
\[
 \begin{split}
 dK_A(a\otimes x)
 &=(-1)^{|a|}da\otimes Kx+a\otimes d_JKx,\\
 K_Ad(a\otimes x)
 &=(-1)^{|a|+1}da\otimes Kx+a\otimes Kd_Jx.
 \end{split}
\]
The terms containing $da$ cancel, leaving
\[
 dK_A+K_Ad=1-\operatorname{const}\operatorname{ev}_0,
 \qquad \operatorname{ev}_0K_A=0.
\]
The second identity follows from $\epsilon_0K=0$.

For a quasi-isomorphism $g:B_0\to A$ of unital differential graded
$S$-algebras, impose the first
endpoint equation by defining
\[
 P(g)=
 \{(b,\omega)\in B_0\oplus(A\otimes J):
                    \operatorname{ev}_0\omega=g(b)\}.
\]
The product and differential preserve this equation because
$\operatorname{ev}_0$ and $g$ are differential graded algebra maps.
The unit is $(1_{B_0},\operatorname{const}(1_A))$.
Thus $P(g)$ is a unital differential graded algebra.
Its degree-$n$ elements have the unique form
\[
 \bigl(b,g(b)\otimes e_0+a\otimes e_1+c\otimes\eta\bigr),
 \qquad b\in B_0^n,\quad a\in A^n,\quad c\in A^{n-1}.
\]

There are unital differential graded algebra maps
\[
 \begin{split}
 j:B_0\longrightarrow P(g),&\qquad
 j(b)=(b,\operatorname{const}g(b)),\\
 p:P(g)\longrightarrow B_0,&\qquad p(b,\omega)=b,\\
 q:P(g)\longrightarrow A,&\qquad
 q(b,\omega)=\operatorname{ev}_1\omega.
 \end{split}
\]
The second endpoint makes $q$ degreewise surjective:
$(0,a\otimes e_1)\in P(g)$ maps to every homogeneous $a\in A$.
The constant inclusion satisfies $pj=1_{B_0}$ and $qj=g$.

Since $\operatorname{ev}_0K_A=0$, the degree-minus-one map
\[
 S(b,\omega)=(0,K_A\omega)
\]
takes values in $P(g)$.
The tensor contraction and the defining endpoint equation give
\[
 \begin{split}
 (dS+Sd)(b,\omega)
 &=(0,\omega-\operatorname{const}\operatorname{ev}_0\omega)\\
 &=(0,\omega-\operatorname{const}g(b))
 =(1-jp)(b,\omega).
 \end{split}
\]
Together with $pj=1$, this proves that $p,j$ induce inverse cohomology maps.
The equation $qj=g$ then proves that $q$ is a quasi-isomorphism.
It provides the surjective backward-arrow replacement needed in
Lemma~\ref{lem:ha-surjective-lifting}.

\begin{proposition}\label{prop:ha-zigzag-transport}
Let $S$ be a commutative unital ring.
Suppose a minimal $A_\infty$-algebra $M$ over $S$ is degreewise
projective and has a quasi-isomorphism to a differential graded
$S$-algebra $A$.
Every finite quasi-isomorphism zigzag of $S$-algebras from $A$ to
$A'$ induces, through any prescribed finite arity, an
$A_\infty$-morphism $M\to A'$ whose linear component is a
quasi-isomorphism.
\end{proposition}

\begin{proof}
For a forward arrow, compose with its strict coalgebra map.
For a backward arrow $g:B_0\to A_0$, replace it by
\[
 B_0\xleftarrow{\ p\ }P(g)\xrightarrow{\ q\ }A_0.
\]
Lift the current morphism $M\to A_0$ through $q$ using
Lemma~\ref{lem:ha-surjective-lifting}, then compose with $p$.
Both maps are quasi-isomorphisms, so the linear component remains
a quasi-isomorphism. Repeat along the finitely many arrows.
The construction keeps the same operations on $M$ throughout.
\end{proof}

For the central algebra $C_r$, return to the original field $k$.
Its cohomology is degreewise free over $k$, so the transport
proposition applies to its transferred minimal structure.

\begin{theorem}\label{thm:ha-formality-dichotomy}
Over a field $k$ of characteristic different from two, the differential graded algebra
$C_r$ is formal over $k$ exactly when $r=1$.
For every $r\geq2$, it is also nonformal over $k[t]$.
\end{theorem}

\begin{proof}
Suppose $r\geq2$ and a formality zigzag exists. Compose the
coordinate isomorphism $M'\to M$ with the transferred morphism
$M\to C_r$. Proposition~\ref{prop:ha-zigzag-transport}, with $N=6$,
then gives a morphism from $M'$ to the zero-differential algebra $H$
through arity six.

Its linear component $\phi$ is a graded vector-space isomorphism.
The arity-two equation gives $\phi(ab)=\phi(a)\phi(b)$.
Thus $\phi$ is a graded algebra automorphism. Postcomposition with
its inverse makes the linear component the identity.
This imposes no constraint on the original action of the zigzag
on $t,u,z$, or on the augmentation.

Let $a=F_3$ and $d=F_5$ be the remaining comparison components.
Their degrees are $-2,-4$, respectively. The evenness of $H$ forces
$F_2=F_4=F_6=0$.

In arity four, equation~\eqref{eq:ha-morphism} reads
\[
 \begin{split}
 a(x_1x_2,x_3,x_4)-a(x_1,x_2x_3,x_4)
 +a(x_1,x_2,x_3x_4)\\
 =x_1a(x_2,x_3,x_4)+a(x_1,x_2,x_3)x_4.
 \end{split}
\]
Hence $\delta a=0$.
In arity six, the source has $m'_6$ and the five alternating
insertions of the binary product into $d$.
The target partitions $(1,5),(3,3),(5,1)$ have positive signs.
Consequently
\begin{equation}\label{eq:ha-formality-comparison}
 m'_6=\delta d+a\smile a.
\end{equation}

The cup square on the right is a boundary. To prove this directly,
write
\[
 E=\delta a,\qquad a_{ijk}=a(x_i,x_j,x_k).
\]
Since the arity of $a$ is three, its self-insertion has three
positive terms. Expansion gives
\[
 \begin{split}
 \delta(a\circ a)(x_1,\ldots,x_6)
 ={}&E(a_{123},x_4,x_5,x_6)
     +E(x_1,a_{234},x_5,x_6)\\
 &+E(x_1,x_2,a_{345},x_6)
     +E(x_1,x_2,x_3,a_{456})\\
 &+a(E(x_1,x_2,x_3,x_4),x_5,x_6)\\
 &-a(x_1,E(x_2,x_3,x_4,x_5),x_6)\\
 &+a(x_1,x_2,E(x_3,x_4,x_5,x_6))
     -2a_{123}a_{456}.
 \end{split}
\]

To check the expansion, the terms containing a multiplication
immediately outside an inner $a$ cancel in pairs. They are
\[
 \begin{gathered}
 a(a_{123}x_4,x_5,x_6),\quad a(x_1a_{234},x_5,x_6),\\
 a(x_1,a_{234}x_5,x_6),\quad a(x_1,x_2a_{345},x_6),\\
 a(x_1,x_2,a_{345}x_6),\quad a(x_1,x_2,x_3a_{456}).
 \end{gathered}
\]
The first and fourth $E$ terms each contribute $a_{123}a_{456}$.
After their subtraction, the remaining terms are precisely the
Hochschild differential of the three self-insertions.

Since $E=0$, the identity reduces to
\[
 \delta(a\circ a)=-2a\smile a.
\]
Equation~\eqref{eq:ha-formality-comparison} would therefore imply
\[
 m'_6=\delta\bigl(d-\tfrac12a\circ a\bigr),
\]
contradicting~\eqref{eq:ha-nonzero-class}.

This excludes all formality zigzags over $k$.
Forgetting the coefficient action turns any $k[t]$-linear zigzag
into a $k$-linear one, proving the second assertion.

When $r=1$, every normal representative lies in $B$.
The inclusion $B\hookrightarrow C_1$ induces the isomorphism
\[
 B\longrightarrow H^*(C_1),\qquad
 H^*(C_1)=B[t]/(t-2uz)\cong B.
\]
Its source has zero differential. It is therefore a unital
quasi-isomorphism proving formality over $k$.
\end{proof}

The class in~\eqref{eq:ha-nonzero-class} obstructs a change of all
coordinates, even one that changes the augmentation.
The augmentation only supplies a linear evaluation proving that this
class is nonzero. The formality contradiction itself permits every
graded algebra automorphism of $H$.

\section{The fourth operation over an arbitrary field}
\label{sec:ha-arbitrary-field}

For this section, let $k$ be a field of arbitrary characteristic.
The preceding six-input argument used $2^{-1}$.
We now compute a fourth-operation obstruction that requires no division.
For the free differential graded algebra
\[
 E_r=k[t]\langle\tau\rangle,\qquad |\tau|=-1,\qquad d\tau=t^r,
\]
this obstruction is nonzero when $r\geq2$.
Four cohomology classes whose adjacent products vanish detect it.
Allowing general even coefficients will also cover $C_r$, its
truncation, and its localization in characteristic two.

Let $B$ be a commutative graded $k$-algebra concentrated in even
degrees, with zero differential. Fix $0\ne u\in B^{-2}$.
Define
\[
 X_r(B,u)=B[t]\langle\tau\rangle/(\tau^2-u),
 \qquad |t|=0,\quad|\tau|=-1,
 \qquad d\tau=t^r,\quad d(B[t])=0,
\]
where $B[t]$ acts centrally.
The free algebra has $B[t]$-basis $1,\tau,\tau^2,\ldots$.

Division by the monic relation in $\tau$ gives the unique expression
\[
 X_r(B,u)=B[t]\oplus\tau B[t].
\]
Its multiplication is
\[
 (a+\tau b)(c+\tau e)=ac+ube+\tau(ae+bc),
 \qquad a,b,c,e\in B[t].
\]
It is associative because it is the product in the displayed quotient
of a free associative algebra.

The graded Leibniz rule preserves the defining relation:
\[
 d(\tau^2-u)=t^r\tau-\tau t^r=0.
\]
Thus the differential descends to the quotient, where
\[
 d(a+\tau b)=t^r b.
\]
This formula gives $d^2=0$ directly.
Multiplication by $t^r$ is injective on $B[t]$, including when $B$
has zero divisors, because it shifts the polynomial coefficients.
Consequently
\[
 H_X:=H^*(X_r(B,u))=B[t]/(t^r),\qquad H_X^{\mathrm{odd}}=0.
\]
The constant coefficient classes identify $B$ with a subalgebra of $H_X$,
so the class of $u$ remains nonzero.

\begin{theorem}\label{thm:ha-common-model-formality}
For every field $k$, every such pair $(B,u)$, and every integer
$r\geq1$, the algebra $X_r(B,u)$ is formal over $k$ exactly when $r=1$.
For every $r\geq1$, it is nonformal over $k[t]$ in the category
of differential graded algebras with central $k[t]$-action.
\end{theorem}

\begin{proof}
Let $i_X:H_X\to B[t]$ choose the representative of $t$-degree less
than $r$. Let $p_X$ take the class of the even component and kill
the odd component. Define
\[
 h_X(a)=\tau\frac{a-i_Xp_X(a)}{t^r},\qquad h_X(\tau b)=0.
\]
The division argument proves
\[
 dh_X+h_Xd=1-i_Xp_X,\qquad
 p_Xi_X=1,\qquad h_X^2=h_Xi_X=p_Xh_X=0.
\]
All these maps are $B$-linear.

The proof of Proposition~\ref{prop:ha-transfer} uses this contraction
identity and the recursion~\eqref{eq:ha-transfer-recursion}.
That recursion and its defect proof use only integer signs and
finite sums, so they apply over the present field.
Define the degree-zero defect $F_X$ by
\[
 i_X(a)i_X(b)-i_X(ab)=t^r\,i_XF_X(a,b).
\]
It is $B$-bilinear, and on powers of $t$ it has the values
\[
 F_X(t^i,t^j)=
 \begin{cases}
 t^{i+j-r},&i+j\geq r,\\
 0,&i+j<r,
 \end{cases}
 \qquad 0\leq i,j<r.
\]

The second comparison component is $f_2^X=-\tau i_XF_X$.
The arity-three outputs are odd and are killed by $p_X,h_X$.
Thus $f_3^X=m_3^X=0$.
At arity four, only the binary partition $(2,2)$ survives,
with its negative sign. Hence
\[
 m_4^X=-uF_X\smile F_X,\qquad \delta m_4^X=0,
\]
where the cocycle identity follows from the arity-five equation.
For example, at $r=2$ this gives
$f_2^X(t,t)=-\tau$ and $m_4^X(t,t,t,t)=-u$.

For general $r\geq2$, put $a=t^{r-1}$.
All three adjacent products of the ordered tuple $(a,t,a,t)$
are zero in $H_X$. Moreover,
\[
 F_X(a,t)=1,\qquad m_4^X(a,t,a,t)=-u.
\]
Choose a linear functional $\ell:B^{-2}\to k$ with $\ell(u)=1$:
extend $u$ to a basis and set $\ell$ to zero on the remaining basis vectors.
Define
\[
 \lambda:H_X^{-2}\longrightarrow k,\qquad
 \lambda\left(\sum_{i=0}^{r-1}b_it^i\right)=\ell(b_0),
 \quad b_i\in B^{-2}.
\]
Then $\lambda(tH_X^{-2})=0$ and $\lambda(u)=1$.

Let $K:H_X^{\otimes3}\to H_X$ be any $k$-linear cochain of degree
$-2$. Its Hochschild differential on this tuple is
\[
 (\delta K)(a,t,a,t)=aK(t,a,t)+K(a,t,a)t.
\]
Every interior term vanishes because its merged adjacent product is zero.
Both remaining terms lie in $tH_X^{-2}$, since $a$ is divisible by $t$.
Consequently
\[
 \lambda(\delta K)(a,t,a,t)=0,\qquad
 \lambda m_4^X(a,t,a,t)=-1\ne0.
\]
Thus $m_4^X$ is not a Hochschild boundary.

The test uses $\lambda$ only as a linear functional.
It requires no augmentation of $B$ or $H_X$, so it also applies
when $u$ is invertible.

A formality zigzag would transport the transferred structure to
strict $H_X$ through arity four, by
Proposition~\ref{prop:ha-zigzag-transport}.
Its linear component is a graded algebra automorphism.
Composing with its inverse makes that component the identity.
Parity forces the second component to vanish, and the arity-four
equation then gives $m_4^X=\delta F_3$.
This contradicts the displayed evaluation.

For $r=1$, the inclusion $B\hookrightarrow X_1(B,u)$ is a
quasi-isomorphism, since it induces $B\cong B[t]/(t)$.

To treat the fixed central action of $k[t]$, let $D$ be an associative
differential graded algebra with $H^{-1}(D)=0$.
Let $s\in D^0$ be a central exact cycle, and choose $x\in D^{-1}$
with $dx=s$. The square is a cycle because
\[
 d(x^2)=sx-xs=0.
\]
Any other primitive is $x+dy$ for some $y\in D^{-2}$.
Indeed, the difference of two primitives is a degree-minus-one cycle,
and the vanishing of $H^{-1}(D)$ makes it a boundary.

The square class is independent of the primitive.
The required identity is
\[
 \begin{split}
 d(-xy+yx+y\,dy)
 &=-sy+x\,dy+(dy)x+ys+(dy)^2\\
 &=(x+dy)^2-x^2.
 \end{split}
\]
The two terms involving $s$ cancel because $s$ is central.
Every differential graded algebra map preserving $s$ preserves this
square class by mapping a primitive and its square to their images.

Along a quasi-isomorphism zigzag preserving the central coefficient
action, the exactness of $s$ and the vanishing of $H^{-1}$ hold at every step.
Both conditions pass in either direction through the induced cohomology
isomorphisms, which preserve the specified class of $s$.
The induced cohomology isomorphisms therefore preserve whether this
square class is zero.

For $X_r(B,u)$, take the fixed scalar $s=t^r$ and its primitive
$x=\tau$. Its square class is $u\ne0$.
In the zero-differential cohomology algebra, $t^r=0$ has primitive zero,
whose square class is zero.
This excludes every central $k[t]$-linear formality zigzag.
The argument uses no inverse of a scalar, so it holds in every characteristic.
\end{proof}

\begin{corollary}\label{prop:ha-witness-k-formality}
Over every field $k$, the algebra $E_r$ is formal over $k$ exactly
when $r=1$. It is nonformal over the central coefficient ring $k[t]$
for every $r\geq1$.
\end{corollary}

\begin{proof}
Take $B=k[u]$ in Theorem~\ref{thm:ha-common-model-formality}.
The unique even--odd decomposition identifies $X_r(k[u],u)$ with $E_r$.
\end{proof}

For $r\geq2$, the fourth operation of $E_r$ has a nonzero
Hochschild class in every characteristic.
For $C_r$ with $2$ invertible, the cochain~\eqref{eq:ha-g-definition}
removes the fourth operation.
The resulting sixth operation has the nonzero class~\eqref{eq:ha-nonzero-class}.

\section{Formality of the nonpositive truncation}
\label{sec:ha-truncation}

Assume again that $k$ is a field of characteristic different from two
in this section and the next.

The positive-degree generator $z$ occurs in the detecting bar chain.
Removing positive degrees changes the formality question. The resulting
algebra admits explicit multiplicative quasi-isomorphisms to cohomology.

For a differential graded algebra $C$, its nonpositive truncation
$\tau_{\leq0}C$ has $C^j$ in degrees $j<0$, the kernel of
$d:C^0\to C^1$ in degree zero, and zero in positive degrees.
The product preserves these subspaces by the Leibniz rule.

In $C_r$, every degree-zero element is a cycle. Set
\[
 A=\tau_{\leq0}C_r,\qquad s=t^r,\qquad v=uz,\qquad B_0=k[t,u].
\]
Here $|v|=0$ and $|u|=-2$.

The nonpositive basis words of $C_r$
are $\tau^n$ and $\tau^{n+2}z$, for $n\geq0$.
Consequently
\[
 \begin{split}
 A&=k[t]\langle\tau\rangle[v]/(v^2)
   =B_0[v]/(v^2)\oplus\tau B_0[v]/(v^2),\\
 D\tau&=s-2v,
 \end{split}
\]
where $v,t$ are central and $u=\tau^2$.
The basis word $\tau^nv$ maps to $\tau^{n+2}z$, proving that
this description is injective as well as surjective.

Multiplication by $s-2v$ is injective on $B_0[v]/(v^2)$.
Indeed,
\[
 (s-2v)(a+vb)=sa+v(sb-2a)=0
\]
first gives $a=0$ and then $b=0$, since multiplication by $t^r$
shifts polynomial coefficients injectively. Thus
\[
 H^*(A)=B_0[v]/(v^2,s-2v)\cong B_0/(s^2).
\]
The last isomorphism eliminates $v=s/2$.

A primitive of $s^2$ in $A$ is $\tau(s+2v)$.
Its square need not vanish as an element, so a single exterior
generator would impose an extra equation. Instead, introduce one
generator for each successive product relation:
\begin{equation}\label{eq:ha-semifree-algebra}
 \begin{gathered}
 P=B_0\langle h_1,h_2,h_3,\ldots\rangle,\qquad |h_n|=1-2n,\\
 dh_1=s^2,\qquad
 dh_n=-\sum_{i+j=n}h_ih_j\quad(n\geq2).
 \end{gathered}
\end{equation}
The sum has positive $i,j$. The coefficients $B_0$ are central,
and their differential is zero. This algebra is semifree over $B_0$:
its underlying algebra is free on the ordered generators, and each
generator's differential uses only earlier ones.

All $h_n$ have odd degree. For $n\geq2$, the Leibniz rule gives
\[
 d^2h_n=-\sum_{i+j=n}(dh_i)h_j+
          \sum_{i+j=n}h_i(dh_j).
\]
The terms with a scalar differential cancel as
$-s^2h_{n-1}+h_{n-1}s^2=0$.
The other terms are two copies of
$\sum_{a+b+c=n}h_ah_bh_c$ with opposite signs.
Together with $d^2h_1=0$, this proves that $P$ is a differential
graded algebra.

Define integers $a_n$ recursively by
\[
 a_1=1,\qquad a_n=-\sum_{i+j=n}a_ia_j\quad(n\geq2).
\]
The recurrence uses integer multiplication and addition only.

There are unital $B_0$-algebra maps
\begin{equation}\label{eq:ha-positive-zigzag}
 B_0/(s^2)\xleftarrow{\ \varepsilon\ }P
          \xrightarrow{\ \Phi\ }A
\end{equation}
given by
\[
 \varepsilon(h_n)=0,\qquad
 \Phi(h_n)=a_n\tau u^{n-1}\bigl(s+(4n-2)v\bigr).
\]
The coefficient maps are the quotient map and the specified action
of $B_0$ on $A$. Every image has degree $1-2n$.

The map $\varepsilon$ respects the differential because $s^2=0$
in its target. For $\Phi$, one computes
\[
 D\Phi(h_n)=a_nu^{n-1}\bigl(s^2+4(n-1)sv\bigr).
\]
For $n=1$ this is $s^2$.
For $i+j=n$, the equality $v^2=0$ gives
\[
 \Phi(h_i)\Phi(h_j)
 =a_ia_j u^{n-1}\bigl(s^2+4(n-1)sv\bigr).
\]
The recurrence for $a_n$ proves $D\Phi(h_n)=\Phi(dh_n)$ on
every generator. Thus both arrows in~\eqref{eq:ha-positive-zigzag}
are differential graded algebra maps.

For $r=2$, the first three generator images are:
\[
 \begin{split}
 \Phi(h_1)&=\tau(t^2+2uz),\\
 \Phi(h_2)&=-\tau u(t^2+6uz),\\
 \Phi(h_3)&=2\tau u^2(t^2+10uz).
 \end{split}
\]
The identity $D\Phi(h_2)=-\Phi(h_1)^2$ recovers the square
calculation at the beginning of the chapter. The higher generators
supply the remaining differential identities.

To prove that these maps are quasi-isomorphisms, compare $P$ with
the differential graded algebra
\[
 K=B_0\otimes\Lambda(h),\qquad |h|=-1,\qquad dh=s^2.
\]
Here $\Lambda(h)$ has basis $1,h$ with $h^2=0$.
There is a differential graded algebra map
\[
 \pi:P\longrightarrow K,\qquad
 \pi(h_1)=h,\qquad\pi(h_n)=0\quad(n\geq2).
\]
For $n=2$, the image of $dh_2$ is $-h^2=0$.
For $n>2$, every product in $dh_n$ contains a generator of index
at least two and also maps to zero.

Give $h_n$ weight $n$, and give $B_0$ weight zero.
Let $F_NP$ be the $B_0$-span of words of weight at most $N$,
and set $F_NP=0$ for $N<0$.
These subspaces form an increasing filtration, meaning a nested
sequence preserved by the differential.
The quadratic part $d_0$ of the differential preserves weight,
while $dh_1=s^2$ lowers it.

Give $h\in K$ weight one and retain weight zero on $B_0$.
The corresponding filtration is
\[
 F_NK=
 \begin{cases}
 0,&N<0,\\
 B_0,&N=0,\\
 K=B_0\oplus B_0h,&N\geq1.
 \end{cases}
\]
The differential $dh=s^2$ lowers weight, so it preserves this
filtration. The map $\pi$ also preserves the filtrations:
it sends $h_1$ to $h$, each $h_n$ with $n\geq2$ to zero,
and fixes $B_0$.

The quotients $\operatorname{gr}_NP=F_NP/F_{N-1}P$ form the
associated graded complex of $P$, with differential $d_0$.
The analogous quotients define $\operatorname{gr}_NK$.
Their differential is zero because $dh=s^2$ lowers weight.
The filtered map $\pi$ induces
$\operatorname{gr}_N\pi:\operatorname{gr}_NP\to
\operatorname{gr}_NK$ on these quotients.

At weight $N\geq1$, the words are indexed by compositions
$N=n_1+\cdots+n_\ell$ into positive integers.
Identify a word with the subset of cuts
\[
 \{n_1,n_1+n_2,\ldots,n_1+\cdots+n_{\ell-1}\}
 \subseteq\{1,\ldots,N-1\}.
\]

Let
\[
 L_N=B_0\otimes_k\Lambda_k(e_1,\ldots,e_{N-1})
\]
as a $B_0$-module.
The exterior factor has basis $e_{j_1}\wedge\cdots\wedge e_{j_q}$
indexed by increasing subsets. Multiplication orders concatenated
symbols with the sign of their permutation, and a repeated symbol
gives zero.
For homogeneous $b\in B_0$, assign
\[
 \bigl|b\,e_{j_1}\wedge\cdots\wedge e_{j_q}\bigr|
 =|b|+q+1-2N.
\]
Send each word to the increasing exterior monomial of its cut
subset, and extend this map $B_0$-linearly.
It is a degree-preserving bijection $\operatorname{gr}_NP\to L_N$.

Splitting the $k$th factor inserts a cut after $k-1$ existing cuts.
The preceding odd factors contribute $(-1)^{k-1}$, and the generator
differential contributes a minus sign.
Thus the differential corresponds to
\[
 d_0=-\omega\wedge(-),\qquad\omega=e_1+\cdots+e_{N-1}.
\]

For $N\geq2$, define the $B_0$-linear map $\iota$ to remove $e_1$
from an increasing monomial containing it and to give zero otherwise.
The exterior relations give
\[
 \iota(\omega\wedge x)+\omega\wedge\iota(x)=x.
\]
Indeed, the $e_1$ summand contributes $x$, while every other
summand appears twice with opposite signs. Thus $c=-\iota$ satisfies
$d_0c+cd_0=1$. In composition notation this contraction is
\[
 c(h_1h_{n_2}\cdots h_{n_\ell})
 =-h_{1+n_2}h_{n_3}\cdots h_{n_\ell},
\]
and it is zero if the first part exceeds one.

Every weight at least two is therefore acyclic.
Weights zero and one contribute $B_0$ and $B_0h_1$.
The associated-graded map of $\pi$ is the identity in those two
weights, and has acyclic source in the remaining weights.

To pass from weights to the full complex, form the cone of $\pi$.
Its degree-$j$ part is $K^j\oplus P^{j+1}$, with differential
\[
 d(y,x)=(d_Ky+\pi x,-d_Px).
\]

Give this cone the filtration
\[
 F_N\operatorname{Cone}(\pi)^j=F_NK^j\oplus F_NP^{j+1}.
\]
Here $F_NK^j=F_NK\cap K^j$ and $F_NP^{j+1}=F_NP\cap P^{j+1}$.
The differential preserves this filtration because $d_K,d_P,\pi$
preserve their filtrations. Taking the two component classes gives
an isomorphism of complexes
\[
 \operatorname{gr}_N\operatorname{Cone}(\pi)
 \cong\operatorname{Cone}(\operatorname{gr}_N\pi).
\]
The displayed cone differential gives the differential on both sides.

At weights zero and one, the cone of the identity contracts by
$(y,x)\mapsto(0,y)$.
At higher weights its target is zero, so the contraction is $-c$
on the shifted source.
Thus every associated graded piece of the cone is acyclic.

For any cycle in this cone, choose its greatest nonzero weight.
Associated-graded acyclicity gives a primitive for that component.
Subtract its differential to obtain a cycle of smaller greatest weight.
Repeating reaches weight zero and then zero after finitely many steps.
Every element has finite weight support, so the argument terminates.

If $y\in K^j$ is a cycle, then $(y,0)$ is a cone cycle.
Acyclicity gives
\[
 (y,0)=d(z,x)=(d_Kz+\pi x,-d_Px).
\]
Thus $x$ is closed and $[y]=H(\pi)[x]$, proving surjectivity.

If $x\in P^j$ is closed and $\pi x=d_Ky$, then $(-y,x)$ is a
cone cycle.
Writing $(-y,x)=d(z,w)$ gives $x=-d_Pw$.
This proves injectivity.

The map $K\to B_0/(s^2)$ sending $h$ to zero is a
quasi-isomorphism. Its odd differential is multiplication by $s^2$,
which is injective, and its even quotient is $B_0/(s^2)$.
Its composite with $\pi$ is $\varepsilon$.
Thus $H^*(P)=B_0/(s^2)$, generated by the coefficient classes.
The map $\Phi$ fixes these classes and maps them to the same
generators of $H^*(A)=B_0/(s^2)$. Hence $\Phi$ is also a
quasi-isomorphism.

\begin{theorem}\label{thm:ha-truncation-formal}
For every $r\geq1$, the nonpositive truncation $\tau_{\leq0}C_r$
is formal over $k[t,u]$. Its cohomology is $k[t,u]/(t^{2r})$,
and~\eqref{eq:ha-positive-zigzag} is an explicit formality zigzag.
\end{theorem}

The same construction and proof work over every commutative
coefficient ring in which $2$ is invertible.
Multiplication by $t^r$ is injective over such a ring because it
shifts polynomial coefficients. The contraction uses only signs,
and the numbers $a_n$ are integers before mapping into the ring.
The inverse of $2$ is required when eliminating $v=s/2$.
For example, in characteristic two the cohomology instead has the
presentation $k[t,u,v]/(v^2,t^r)$, so that elimination is unavailable.

\section{Localizing at an even cycle}
\label{sec:ha-localization}

The inclusion $A\hookrightarrow C_r$ omits $z$ but contains $v=uz$.
After $u$ becomes invertible, the equation $z=u^{-1}v$ recovers
the omitted generator.

Let $C$ be a graded algebra with a central element $u$ of degree $-2$.
Its localization consists of fractions $u^{-n}a$, with $n\geq0$.
Equality means
\[
 u^{-n}a=u^{-m}b
 \quad\Longleftrightarrow\quad
 u^\ell(u^ma-u^nb)=0
 \quad\text{for some }\ell\geq0.
\]
Addition uses a common denominator, and
\[
 (u^{-n}a)(u^{-m}b)=u^{-(n+m)}ab.
\]
A homogeneous fraction has degree $|a|+2n$.

If $C$ has differential $D$ with $Du=0$, set
$D(u^{-n}a)=u^{-n}Da$.
These operations respect fraction equality: multiplying an equality
by a numerator uses centrality, and differentiating it uses $Du=0$
and the even degree of $u$.
The differential still squares to zero and satisfies the Leibniz rule,
as follows by applying those identities to the numerators.
Only finite sums and finite denominators occur.

The displayed equation for $z$ gives an isomorphism
\[
 A[u^{-1}]\longrightarrow C_r[u^{-1}].
\]
Its inverse sends $z$ to $u^{-1}v$ and preserves $t,\tau,u$.
The relation $z^2=0$ and the formula for $D\tau$ are preserved
because $v^2=0$ and $uz=v$.

Localization commutes with cohomology.
Sending a fraction of a closed representative to its class defines
an algebra map
\[
 H^*(C)[u^{-1}]\longrightarrow H^*(C[u^{-1}]).
\]
A boundary in the numerator maps to a boundary, and the fraction
relation makes the map independent of the representative fraction.
If $u^{-n}a$ is closed, then $u^mDa=0$ for some $m\geq0$.
Thus
\[
 u^{-n}a=u^{-(n+m)}u^ma
\]
has a closed numerator, proving surjectivity.

For injectivity, let $a$ be closed and suppose the image of
$u^{-j}[a]$ is a boundary.
Multiplying by the invertible cycle $u^j$ gives
$a=D(u^{-n}b)$ for some $n\geq0$.
Fraction equality gives
\[
 u^m(u^na-Db)=0
\]
for some $m\geq0$.
Hence $u^{m+n}a=D(u^mb)$ before localization, so
$u^{-j}[a]=0$ in $H^*(C)[u^{-1}]$.
The same argument applies to a complex of graded $k[u]$-modules.

Applying this to the two quasi-isomorphisms in
\eqref{eq:ha-positive-zigzag} gives
\[
 k[t,u,u^{-1}]/(t^{2r})
 \xleftarrow{\ \simeq\ }P[u^{-1}]
 \xrightarrow{\ \simeq\ }C_r[u^{-1}].
\]

This is a zigzag over $k[t,u,u^{-1}]$.
In the cohomology of its right-hand algebra,
\[
 z=\frac{t^r}{2u},\qquad z^2=0
 \quad\Longleftrightarrow\quad t^{2r}=0.
\]
Thus the localized full algebra is formal over the displayed
coefficient ring for every $r\geq1$.
These maps specify the action of $t,u,u^{-1}$.
They do not assert formality with an additional fixed action of $z$.

The nonformality obstruction uses both cohomology relations
$t^r=2uz$ and $z^2=0$.
Truncation removes $z$ from the algebra.
After localization, its cohomology class satisfies $z=t^r/(2u)$.
The displayed zigzags prove formality for these changed algebras.

\section{Characteristic two and the fixed coefficient action}
\label{sec:ha-characteristic-two}

Let $k$ now be a field of characteristic two, and retain the defining
formulas for $C_r$. Its differential reduces to $D\tau=t^r$.
The relation coupling $u$ and $z$ in the differential has disappeared.
The full algebra, its nonpositive truncation, and its localization
therefore belong to the family $X_r(B,u)$.

Set $A_r=\tau_{\leq0}C_r$, and define the even coefficient algebras
\[
 B_C=k[u,z]/(z^2),\qquad B_A=k[u,v]/(v^2),\qquad
 B_L=k[u,u^{-1},z]/(z^2),
\]
with $|u|=-2$, $|z|=2$, and $|v|=0$.
All have zero differential, and $u$ is nonzero in each.
Their presentations give
\[
 \begin{split}
 C_r&=X_r(B_C,u),\\
 A_r&=X_r(B_A,u),\\
 C_r[u^{-1}]&=X_r(B_L,u).
 \end{split}
\]
The first identity follows by separating the even and odd powers
of $\tau$.

For the second identity, the nonpositive words are $\tau^n$ and
$\tau^{n+2}z$, with $n\geq0$.
Fix $t,\tau$, and send $v$ to $uz$, so $\tau^nv$ maps to $\tau^{n+2}z$.
These are bijections of the displayed bases, preserving products and
the differential $D\tau=t^r$.
For the third identity, adjoining the central inverse $u^{-1}$
retains the decomposition into even and odd powers of $\tau$.
The equality $\tau^{-1}=\tau u^{-1}$ introduces no further generators.

The common cohomology calculation gives
\[
 \begin{split}
 H^*(C_r)&=k[t,u,z]/(t^r,z^2),\\
 H^*(A_r)&=k[t,u,v]/(t^r,v^2),\\
 H^*(C_r[u^{-1}])&=k[t,u,u^{-1},z]/(t^r,z^2).
 \end{split}
\]
In particular, the class of $u$ survives in all three algebras.

\begin{theorem}\label{thm:ha-characteristic-two}
Over a field of characteristic two, each of $C_r$, $A_r$, and
$C_r[u^{-1}]$ is formal over $k$ exactly when $r=1$.
Each is nonformal over the central coefficient ring $k[t]$ for
every $r\geq1$.
Consequently $A_r$ is nonformal over $k[t,u]$, and
$C_r[u^{-1}]$ is nonformal over $k[t,u,u^{-1}]$, with their central
coefficient actions fixed.
\end{theorem}

\begin{proof}
The three presentations satisfy the hypotheses of
Theorem~\ref{thm:ha-common-model-formality}.
That theorem gives both the formality classifications over $k$ and
the nonformality over central $k[t]$.
A zigzag fixing either of the larger coefficient rings would also
fix $k[t]$, so it is excluded.
\end{proof}

For $r=1$, the inclusions of $B_C,B_A,B_L$ into their respective
algebras give the asserted formality over $k$.
Their source cohomology algebras have $t=0$, while the target
algebras contain the nonzero element $t=D\tau$.
Thus these inclusions do not preserve the chain-level action of $k[t]$.
The square class in Theorem~\ref{thm:ha-common-model-formality}
excludes every comparison that does preserve it.

For every field, $C_r$ is therefore formal over $k$ exactly when
$r=1$.
The algebras $A_r$ and $C_r[u^{-1}]$ are formal over $k$ exactly
when $2\ne0$ in $k$ or $r=1$.
Their formality over the respective central coefficient rings
$k[t,u]$ and $k[t,u,u^{-1}]$ holds exactly when $2\ne0$ in $k$.

\section{Cohomology under reduction of coefficients}
\label{sec:ha-integral-reduction}

Cohomology can commute with every extension of coefficients while
formality changes under reduction.
An algebra over the integers with the same presentation as the
nonpositive truncation exhibits this distinction.
Fix $r\geq2$, and put
\[
 T_{\mathbb Z}=(\mathbb Z[u,v]/(v^2))[t],
 \qquad |u|=-2,\quad |v|=|t|=0.
\]
Define
\[
 A_{\mathbb Z}=T_{\mathbb Z}\oplus\tau T_{\mathbb Z},
 \qquad |\tau|=-1,\qquad \tau^2=u,
\]
with $\tau$ commuting with $T_{\mathbb Z}$.

The multiplication and differential are
\[
 \begin{split}
 (a+\tau b)(c+\tau e)&=ac+ube+\tau(ae+bc),\\
 d(a+\tau b)&=(t^r-2v)b.
 \end{split}
\]
Equivalently, the underlying algebra is
$T_{\mathbb Z}\langle\tau\rangle/(\tau^2-u)$.
Its differential kills $T_{\mathbb Z}$ and sends $\tau$ to $t^r-2v$.
The latter element commutes with $\tau$, so
$d(\tau^2-u)=(t^r-2v)\tau-\tau(t^r-2v)=0$.
Thus the graded Leibniz rule descends to this quotient, and the
displayed differential gives $d^2=0$.

For a commutative ring $S$, set
$A_S=A_{\mathbb Z}\otimes_{\mathbb Z}S$.
The tensor product has the same presentation over $S$, with the
integer $2$ interpreted in $S$.
Indeed, the words $u^av^jt^i$ and $\tau u^av^jt^i$, with
$a,i\geq0$ and $j=0,1$, give a free abelian basis before extending
coefficients. Extending this basis and the displayed structure
formulas gives precisely that presentation over $S$.

The polynomial $t^r-2v$ is monic in $t$ over $S[u,v]/(v^2)$.
Its multiplication is therefore injective, and division gives
\[
 H^*(A_S)=S[u,v,t]/(v^2,t^r-2v).
\]
This cohomology has no odd-degree part.
Its normal representatives form the free $S$-module on
\[
 u^av^jt^i,\qquad a\geq0,\quad j=0,1,\quad 0\leq i<r.
\]
The degree of this monomial is $-2a$.
In particular, $H^{-2a}(A_{\mathbb Z})$ is free abelian of rank
$2r$ for every $a\geq0$, and all other degrees vanish.

In degree zero, the relation $t^r=2v$ gives $t^{2r}=4v^2=0$.
Thus the scalar map factors as
\[
 \pi_0:\mathbb Z[t]/(t^{2r})\longrightarrow H^0(A_{\mathbb Z}).
\]
Use the ordered bases
\[
 (1,t,\ldots,t^{2r-1})
 \quad\text{and}\quad
 (1,t,\ldots,t^{r-1},v,vt,\ldots,vt^{r-1})
\]
for its source and target.
For $0\leq i<r$, the first $r$ basis vectors map to $t^i$,
and the remaining vectors satisfy $\pi_0(t^{r+i})=2vt^i$.
Writing $I_r$ for the $r$-by-$r$ identity matrix, the matrix of $\pi_0$ is
\[
 \begin{pmatrix}I_r&0\\0&2I_r\end{pmatrix}.
\]

Multiplication by $2$ is injective on $\mathbb Z$, so $\pi_0$ is injective.
Its cokernel is $(\mathbb Z/2\mathbb Z)^r$.
Over the rational field $\mathbb Q$, the same matrix is invertible.
Let $\mathbb F_2=\mathbb Z/(2)$.
Over this field, the last block is zero.
The scalar comparison over $\mathbb F_2$ therefore has the
$r$-dimensional kernel
\[
 \operatorname{span}_{\mathbb F_2}
 \{t^r,t^{r+1},\ldots,t^{2r-1}\}.
\]

Consequently the scalar kernels over $\mathbb Z$ and $\mathbb Q$
are $(t^{2r})$, while the scalar kernel over $\mathbb F_2$ is $(t^r)$.
These are also the ideals that annihilate the full cohomology:
a scalar kills every class exactly when it kills the unit.
The same normal cohomology basis survives the reduction, while the
scalar comparison loses injectivity.

\begin{corollary}\label{cor:ha-formality-under-reduction}
For every commutative ring $S$, the canonical map
\[
 H^*(A_{\mathbb Z})\otimes_{\mathbb Z}S
 \longrightarrow H^*(A_{\mathbb Z}\otimes_{\mathbb Z}S)
\]
is an isomorphism.
The algebra $A_{\mathbb Z}\otimes_{\mathbb Z}\mathbb Q$ is formal
over $\mathbb Q$, whereas
$A_{\mathbb Z}\otimes_{\mathbb Z}\mathbb F_2$ is nonformal over
the two-element field $\mathbb F_2$.
\end{corollary}

\begin{proof}
The canonical map sends $[x]\otimes s$ to $[x\otimes s]$.
It is multiplicative by the tensor-product multiplication.
It sends each displayed normal-basis class to the corresponding
normal-basis class over $S$.
These are free-module bases on both sides, so the map is an isomorphism.
This proof imposes no flatness condition on $S$.

Over $\mathbb Q$, the presentation of $A_S$ is the nonpositive
truncation model in Theorem~\ref{thm:ha-truncation-formal}.
Its explicit zigzag proves formality.
Over $\mathbb F_2$, the differential is $d\tau=t^r$, so
\[
 A_{\mathbb F_2}
 =X_r\bigl(\mathbb F_2[u,v]/(v^2),u\bigr).
\]
Theorem~\ref{thm:ha-common-model-formality} proves nonformality because
$u\ne0$ and $r\geq2$.
\end{proof}

Define $P_{\mathbb Z}$ and
$\Phi_{\mathbb Z}:P_{\mathbb Z}\to A_{\mathbb Z}$ by the integer
formulas in~\eqref{eq:ha-semifree-algebra} and~\eqref{eq:ha-positive-zigzag},
with $B_0=\mathbb Z[t,u]$ and $s=t^r$.
The recurrence makes $\Phi_{\mathbb Z}$ a chain map without dividing by two.
The $B_0$-linear weight contraction and the finite cone argument use
only integer signs, while multiplication by $t^{2r}$ is injective.
They therefore give
\[
 H^*(P_{\mathbb Z})=\mathbb Z[t,u]/(t^{2r}).
\]

The map $\Phi_{\mathbb Z}$ fixes $t,u$.
Under the normal-basis identifications, its cohomology map in degree
$-2a$ is $u^a\pi_0$, with matrix $\operatorname{diag}(I_r,2I_r)$.
The same contraction computes $H^*(P_{\mathbb Z}\otimes S)$
for $S=\mathbb Q$ and $S=\mathbb F_2$.
Thus this cohomology comparison is an isomorphism over $\mathbb Q$
and is singular over $\mathbb F_2$, although the cohomology of
$A_{\mathbb Z}$ commutes with both extensions of coefficients.
The fourth-operation obstruction excludes every replacement
formality zigzag over $\mathbb F_2$.

\section{The exact coefficient-ring boundary}
\label{sec:ha-ring-formality}

The coefficient two in the relation $t^r=2v$ determines both the
fourth-operation detector and the scalar comparison matrix.
We now replace it by an arbitrary scalar and determine which
property of that scalar permits formality.
Let $S$ be a commutative unital ring, let $\kappa\in S$, and let
$r\geq1$ be an integer.
All algebras and maps in this section have central $S$-action,
and all tensor products and polynomials are algebraic.
An element is a unit if it has a multiplicative inverse, and
$S^\times$ denotes the set of units.
The zero ring is allowed, with $0=1$ and its unique element a unit.

Put
\[
 B_S=S[u,v]/(v^2),\qquad T_S=B_S[t],\qquad f_\kappa=t^r-\kappa v,
 \qquad |u|=-2,\quad |v|=|t|=0.
\]
Define
\[
 A_\kappa=A(S,\kappa,r)=T_S\oplus\tau T_S,
 \qquad |\tau|=-1,\qquad\tau^2=u,
\]
where $\tau$ commutes with $T_S$, and set
\[
 \begin{split}
 (a+\tau b)(c+\tau e)&=ac+ube+\tau(ae+bc),\\
 d(a+\tau b)&=f_\kappa b.
 \end{split}
\]
The quotient $T_S\langle\tau\rangle/(\tau^2-u)$ gives this
associative product and unit.
There is no imposed relation $2\tau^2=0$.
The degree-one differential descends by the graded Leibniz rule,
since $d(\tau^2-u)=f_\kappa\tau-\tau f_\kappa=0$,
and its displayed formula squares to zero.

Multiplication by $f_\kappa$ is injective on $T_S$.
Indeed, multiplying a nonzero polynomial by this monic polynomial
retains a nonzero leading coefficient, even when $S$ has zero divisors.
Consequently
\[
 H_\kappa:=H^*(A_\kappa)
 =S[u,v,t]/(v^2,t^r-\kappa v),\qquad H_\kappa^{\mathrm{odd}}=0.
\]
Monic division gives the $S$-basis
\[
 u^av^jt^i,\qquad a\geq0,\quad j=0,1,\quad 0\leq i<r.
\]
The monomial has degree $-2a$, so each $H_\kappa^{-2a}$ is free
on the $2r$ displayed generators, and all other degrees vanish.
For the zero ring these statements mean equalities of zero modules.

\begin{theorem}\label{thm:ha-ring-truncation-formality}
For every commutative unital ring $S$, scalar $\kappa\in S$, and
integer $r\geq1$,
\[
 A(S,\kappa,r)\text{ is formal over }S
 \quad\Longleftrightarrow\quad
 r=1\ \text{or}\ \kappa\in S^\times.
\]
Here formality is in the category of unital associative differential
graded algebras with central $S$-action, and means a finite zigzag
of unital $S$-linear quasi-isomorphisms to cohomology with zero differential.
When $r=1$, a formality map preserves the action of $B_S$.
When $\kappa$ is a unit, a formality zigzag preserves the action of $S[t,u]$.
\end{theorem}

Passing a hypothetical zigzag to $S/\kappa S$ by ordinary tensor
product would not prove the negative implication.
Even free cohomology does not ensure that a quasi-isomorphism
survives such a tensor product.
For an explicit example, let $L$ be the complex
\[
 L^{-2}=\mathbb Z\xrightarrow{\ 2\ }L^{-1}=\mathbb Z
 \xrightarrow{\ \bmod 2\ }L^0=\mathbb Z/2,
\]
with zero terms in other degrees.
Its first map is injective, the kernel of its second map is $2\mathbb Z$,
and its second map is surjective, so $L$ is acyclic.

Form the differential graded algebra $\mathbb Z\oplus L$ by
putting all products of two elements of $L$ equal to zero and
using the usual scalar action of $\mathbb Z$.
Explicitly,
\[
 (s,x)(t,y)=(st,sy+tx),\qquad d(s,x)=(0,dx).
\]
For factors $(s,x),(t,y),(w,z)$, both triple products have scalar
component $stw$ and $L$-component $stz+swy+twx$.
Thus the product is associative.
The unit is $(1,0)$, and the differential satisfies the Leibniz rule:
scalar factors have degree zero, and every product within $L$ is zero.
The projection $\mathbb Z\oplus L\to\mathbb Z$ is therefore a
unital quasi-isomorphism with free cohomology.

After tensoring with $\mathbb F_2$, the ideal complex becomes
\[
 \mathbb F_2\xrightarrow{\ 0\ }\mathbb F_2
 \xrightarrow{\ 1\ }\mathbb F_2.
\]
It has $H^{-2}=\mathbb F_2$, so the tensored projection is not
a quasi-isomorphism.
The acyclic kernel $L$ also has no contraction: in degree zero a
contraction would supply a section $\mathbb Z/2\to\mathbb Z$
of the quotient map, but every such homomorphism is zero.
The projective-source transport proof works over $S$ throughout,
so it avoids both of these failed implications.

\begin{proof}[Proof of Theorem~\ref{thm:ha-ring-truncation-formality}]
For the zero ring, all the displayed algebras are zero and $\kappa$
is a unit, so the conclusion holds.
Assume henceforth that $S$ is nonzero.
Let $i_\kappa:H_\kappa\to T_S$ select the normal representative,
and let $p_\kappa:A_\kappa\to H_\kappa$ take the class of the
even component and kill the odd component.
Set
\[
 h_\kappa(a)=\tau\frac{a-i_\kappa p_\kappa(a)}{f_\kappa},
 \qquad h_\kappa(\tau b)=0.
\]
Monic division gives
\[
 dh_\kappa+h_\kappa d=1-i_\kappa p_\kappa,\qquad
 p_\kappa i_\kappa=1,\qquad
 h_\kappa^2=h_\kappa i_\kappa=p_\kappa h_\kappa=0.
\]

These maps give a transferred minimal structure over $S$ by the
recursion~\eqref{eq:ha-transfer-recursion}.
The defect calculation in Proposition~\ref{prop:ha-transfer} uses
only this contraction and the signed coalgebra identities already
proved over $S$ in Section~\ref{sec:ha-transport}.
Its final injectivity step also holds: $I=si_\kappa s^{-1}$ has
left inverse $P=sp_\kappa s^{-1}$, so
$P^{\otimes n}I^{\otimes n}=1$ over $S$.
Taking the greatest nonzero input length again makes the coalgebra
map injective and forces the transferred coderivation to square to zero.
Thus this explicit contraction supplies the minimal model.

Define the degree-zero, $B_S$-bilinear defect $F_\kappa$ by
\[
 i_\kappa(a)i_\kappa(b)-i_\kappa(ab)
 =f_\kappa\,i_\kappa F_\kappa(a,b).
\]
For $0\leq i,j<r$, it has the formula
\[
 F_\kappa(t^i,t^j)=
 \begin{cases}
 t^{i+j-r},&i+j\geq r,\\
 0,&i+j<r.
 \end{cases}
\]
When $i+j\geq r$, the identity
$t^{i+j}=f_\kappa t^{i+j-r}+\kappa vt^{i+j-r}$ is the division
formula, because $i+j-r<r$.
When $i+j<r$, the product already has normal degree.
The $B_S$-linearity of division proves the asserted bilinearity.

For an even graded $S$-algebra, use $S$-linear Hochschild cochains
and tensor products over $S$ in the differential and cup formulas
of Section~\ref{sec:ha-hochschild}.
The proofs of $\delta^2=0$ and the cup differential use only signed
multilinear expansions, so they hold over $S$.
The transfer gives $f_2=-\tau i_\kappa F_\kappa$.
Every arity-three binary term has odd output and is killed by
$p_\kappa,h_\kappa$, so $f_3=m_3=0$.

Only the partition $(2,2)$ contributes at arity four.
Its sign is negative, and the even input degrees introduce no
additional tensor sign, so
\[
 m_4=-uF_\kappa\smile F_\kappa,\qquad \delta m_4=0.
\]
The second identity is the arity-five relation, with its common
sign removed.
In the smallest case $r=2$, the formula gives
$m_4(t,t,t,t)=-u$, independently of $\kappa$.
The distinction between scalars will arise from the Hochschild boundaries.

Suppose $r\geq2$ and $A_\kappa$ is formal over $S$.
The normal basis makes $H_\kappa$ degreewise free and therefore
projective over $S$.
Proposition~\ref{prop:ha-zigzag-transport} transports its fixed
minimal model through the arbitrary formality zigzag to strict
$H_\kappa$, using only arities through four.
The linear component is an $S$-linear graded algebra automorphism
by the arity-two equation.
Composing with its inverse makes this component the identity,
without fixing the original coordinates $t,u,v$ along the zigzag.

Evenness forces the second comparison component to vanish.
The arity-four equation therefore gives
\[
 m_4=\delta K
\]
for an $S$-linear degree-minus-two cochain
$K:H_\kappa^{\otimes_S3}\to H_\kappa$.
Put $a=t^{r-1}$.
The three adjacent products of $(a,t,a,t)$ are all $\kappa v$, and
\[
 F_\kappa(a,t)=1,\qquad m_4(a,t,a,t)=-u.
\]

The assignment $t,v\mapsto0$ defines a graded algebra map
$H_\kappa\to S[u]$, since it sends both defining relations to zero.
On degree minus two, compose it with $bu\mapsto b$ for $b\in S$
to obtain
\[
 \lambda:H_\kappa^{-2}\longrightarrow S.
\]
Thus $\lambda(u)=1$ and $\lambda(tH_\kappa^{-2})=0$.
In terms of the normal basis, $\lambda$ extracts the coefficient
of $u$ with no $v$ or $t$ factor.

The Hochschild differential on the chosen tuple is
\[
 \begin{split}
 (\delta K)(a,t,a,t)
 ={}&aK(t,a,t)+K(a,t,a)t\\
 &+\kappa\bigl(-K(v,a,t)+K(a,v,t)-K(a,t,v)\bigr).
 \end{split}
\]
The three interior products equal $\kappa v$, and $S$-linearity
extracts $\kappa$ with the displayed alternating signs.
Because $r\geq2$, the element $a$ is divisible by $t$, so
$\lambda$ kills both exterior terms.
Consequently $m_4=\delta K$ implies
\[
 -1=\kappa b,\qquad
 b=\lambda\bigl(-K(v,a,t)+K(a,v,t)-K(a,t,v)\bigr)\in S.
\]
Thus $-b$ is an inverse of $\kappa$.

Equivalently, the functional
\[
 P\longmapsto\lambda(P(a,t,a,t))\bmod\kappa S
 \quad\text{on}\quad
 C_S^{4,-2}(H_\kappa,H_\kappa)
 =\operatorname{Hom}_S^{-2}(H_\kappa^{\otimes_S4},H_\kappa)
\]
kills every Hochschild boundary and sends $m_4$ to $-1$ in
$S/\kappa S$.
If $\kappa$ is not a unit, then $1\notin\kappa S$, so this quotient
and this value are nonzero.
Only the explicit cochain equation is evaluated modulo $\kappa S$.
No arrow of the original zigzag is tensored with this quotient.

For the first positive case, $r=1$, normal forms identify
\[
 H_\kappa=B_S[t]/(t-\kappa v)\cong B_S.
\]
The inclusion $B_S\hookrightarrow A_\kappa$ is a unital
quasi-isomorphism and preserves the displayed $B_S$-action.
The cohomological action of $t$ on its source is multiplication by
$\kappa v$, while the target retains $t$ as a polynomial element.
Thus this inclusion does not fix the chain-level action of $t$.

To treat a unit $\kappa$, construct a comparison for every scalar
$\kappa$ and compute when it is a quasi-isomorphism.
Put $R_S=S[t,u]$ and $s=t^r$.
Since $s^2=\kappa^2v^2=0$ in $H_\kappa$, the coefficient map
$R_S\to H_\kappa$ factors through $R_S/(s^2)$.
A primitive for $s^2$ in $A_\kappa$ is
\[
 w=\tau(s+\kappa v),\qquad dw=(s-\kappa v)(s+\kappa v)=s^2.
\]
Its square is $u(s^2+2\kappa sv)$, a nonzero element when $S\ne0$,
because its coefficient of $ut^{2r}$ is one.

An exterior generator $h$ with $h^2=0$ therefore cannot map to this
primitive, and injectivity of $d$ on $\tau T_S$ makes the primitive unique.
Instead, use the free associative algebra
\[
 P_S=R_S\langle h_1,h_2,\ldots\rangle,\qquad |h_n|=1-2n,
\]
with central coefficient ring $R_S$, zero differential on $R_S$, and
\[
 dh_1=s^2,\qquad
 dh_n=-\sum_{\substack{i+j=n\\i,j\geq1}}h_ih_j\quad(n\geq2).
\]
Each differential is a finite expression in earlier generators.
For $n\geq2$, the Leibniz rule gives
\[
 d^2h_n=-\sum_{i+j=n}(dh_i)h_j+\sum_{i+j=n}h_i(dh_j).
\]
The scalar terms cancel as $-s^2h_{n-1}+h_{n-1}s^2=0$,
and the other terms are the two opposite sums of $h_ah_bh_c$ over
$a+b+c=n$ with positive indices.
Together with $d^2h_1=0$, this proves $d^2=0$.

Use the integers $a_n$ defined by
\[
 a_1=1,\qquad a_n=-\sum_{i+j=n}a_ia_j\quad(n\geq2).
\]
They give unital graded $R_S$-algebra maps
\[
 R_S/(s^2)\xleftarrow{\ \varepsilon\ }P_S
 \xrightarrow{\ \Phi_\kappa\ }A_\kappa
\]
by
\[
 \varepsilon(h_n)=0,\qquad
 \Phi_\kappa(h_n)=a_n\tau u^{n-1}\bigl(s+(2n-1)\kappa v\bigr).
\]
The left coefficient map is the quotient, and the right coefficient
map is the specified action on $A_\kappa$.
The image of $h_n$ has degree $1-2n$.
The left arrow commutes with the differential because its target
has $s^2=0$.

For the right arrow, $v^2=0$ gives
\[
 \begin{split}
 d\Phi_\kappa(h_n)
 &=a_nu^{n-1}(s-\kappa v)(s+(2n-1)\kappa v)\\
 &=a_nu^{n-1}\bigl(s^2+2(n-1)\kappa sv\bigr).
 \end{split}
\]
For $n=1$ this is $s^2$.
For positive $i,j$ with $i+j=n$, multiplication gives
\[
 \begin{split}
 \Phi_\kappa(h_i)\Phi_\kappa(h_j)
 &=a_ia_j u^{n-1}
   (s+(2i-1)\kappa v)(s+(2j-1)\kappa v)\\
 &=a_ia_j u^{n-1}\bigl(s^2+2(n-1)\kappa sv\bigr).
 \end{split}
\]
The recurrence proves $d\Phi_\kappa(h_n)=\Phi_\kappa(dh_n)$
on every generator, for every scalar $\kappa$.
In particular, the first two images satisfy
$\Phi_\kappa(h_1)=w$ and $d\Phi_\kappa(h_2)=-w^2$.

To compute the cohomology of $P_S$, put
\[
 K_S=R_S\oplus R_Sh,\qquad |h|=-1,\qquad h^2=0,\qquad dh=s^2.
\]
The square-zero relation is imposed in every characteristic.
The graded Leibniz rule preserves it because $s^2h-hs^2=0$.
The assignment $h_1\mapsto h$, $h_n\mapsto0$ for $n\geq2$
defines a unital $R_S$-algebra map $\pi:P_S\to K_S$.
It respects the differential: the image of $dh_2$ is $-h^2=0$,
and each product in $dh_n$ for $n>2$ contains a generator with index
at least two.

Give $h_n$ weight $n$, give $h$ weight one, and give $R_S$ weight zero.
The filtrations by weight at most $N$ and their cone filtration
are those of Section~\ref{sec:ha-truncation}, now over $R_S$.
At weight $N\geq1$, compositions of $N$ give the cut-module description
\[
 \operatorname{gr}_NP_S
 \cong R_S\otimes_S\Lambda_S(e_1,\ldots,e_{N-1}),
 \qquad |b\,e_{j_1}\wedge\cdots\wedge e_{j_q}|=|b|+q+1-2N
\]
for homogeneous $b\in R_S$.
A split in the $k$th factor contributes $(-1)^{k-1}$ from the
preceding odd factors and a minus sign from the generator differential.
Hence the differential on this module is $-\omega\wedge(-)$,
where $\omega=e_1+\cdots+e_{N-1}$.

For $N\geq2$, let $\iota$ be the $R_S$-linear map removing $e_1$
from an increasing exterior monomial containing it and giving zero otherwise.
The $e_1$ term and cancellation of all other terms give
\[
 \iota(\omega\wedge x)+\omega\wedge\iota(x)=x.
\]
Thus $-\iota$ contracts every weight at least two over $S$.
In weights zero and one, $\operatorname{gr}\pi$ is the identity.
The associated graded cone is therefore acyclic, using the explicit
identity-cone and shifted-source contractions already derived in
Section~\ref{sec:ha-truncation}.

For each cycle in the cone, subtract a differential lifting a primitive
of its greatest nonzero weight component.
The greatest weight strictly decreases, so the nonnegative filtration
and finite support make this process terminate.
The cone is acyclic, and the two cone equations in
Section~\ref{sec:ha-truncation} prove that $\pi$ is a quasi-isomorphism.
Finally, multiplication by $s^2=t^{2r}$ is injective on $R_S$, so
\[
 H^*(P_S)\cong H^*(K_S)\cong R_S/(s^2).
\]
The composite $P_S\to K_S\to R_S/(s^2)$ is $\varepsilon$,
which is therefore a quasi-isomorphism for every $\kappa$.

The right arrow fixes $t,u$, so its cohomology map is the scalar map
\[
 H(\Phi_\kappa):R_S/(s^2)\longrightarrow H_\kappa.
\]
In degree $-2a$, for every $a\geq0$, use the ordered bases
\[
 \begin{split}
 &u^a(1,t,\ldots,t^{2r-1}),\\
 &u^a(1,t,\ldots,t^{r-1},v,vt,\ldots,vt^{r-1})
 \end{split}
\]
for its source and target.
For $0\leq i<r$, the first $r$ vectors map to themselves,
while $t^{r+i}$ maps to $\kappa vt^i$.
Thus the matrix in every such degree is
\[
 \operatorname{diag}(I_r,\kappa I_r).
\]

Write $\operatorname{Ann}_S(\kappa)=\{b\in S:\kappa b=0\}$.
The first block has zero kernel and cokernel, while each coordinate
of the second block is multiplication by $\kappa$.
Consequently, in every degree $-2a$,
\[
 \ker H(\Phi_\kappa)\cong\operatorname{Ann}_S(\kappa)^r,
 \qquad
 \operatorname{coker}H(\Phi_\kappa)\cong(S/\kappa S)^r.
\]
All other degrees have zero source and target.
This map is an isomorphism exactly when $\kappa$ is a unit:
a unit makes the second block invertible, while surjectivity gives
$1\in\kappa S$ and hence an inverse.
When $\kappa$ is a unit, both arrows in the displayed zigzag are
therefore quasi-isomorphisms over $S[t,u]$, completing the proof.
\end{proof}

When $r=1$ and $\kappa$ is a nonunit, $\Phi_\kappa$ still fails
to be a quasi-isomorphism, but the inclusion of $B_S$ proves formality.
For $r\geq2$, the functional modulo $\kappa S$ excludes every
replacement formality zigzag when this scalar map is not invertible.
Thus the matrix tests the displayed comparison, while the fourth
operation establishes the full nonexistence statement.

The family with fixed coefficient two is the specialization
\[
 A_{S,r}=A(S,2\cdot1_S,r).
\]
It is formal over $S$ exactly when $r=1$ or $2$ is a unit in $S$.
For $r\geq2$, the integral algebra $A_{\mathbb Z}$ of
Section~\ref{sec:ha-integral-reduction} is nonformal over $\mathbb Z$,
and its scalar extension to $S$ is formal exactly when $2$ becomes a unit.
This scalar extension has $\kappa=2\cdot1_S$.
An arbitrary scalar $\kappa$ defines a larger family and is not
in general obtained by extending that fixed integral algebra.

For $r\geq2$, $A(\mathbb Z/4,2,r)$ is nonformal, since its detector
has value $-1=3$ modulo the ideal $2S=\{0,2\}$.
The algebra $A(\mathbb Z/4,3,r)$ is formal for every $r$, since
$3$ is a unit even though the ring contains a nonzero nilpotent.
Over $\mathbb Z[1/2]$, the fixed-coefficient-two family is formal
for every $r$.
Over a field, the condition on $\kappa$ is simply $\kappa\ne0$,
with the additional formal case $r=1$ for every $\kappa$.
In particular, $A(\mathbb F_2,1,r)$ is formal for every $r$,
whereas $A(\mathbb F_2,0,r)$ is nonformal for $r\geq2$.

In any nonzero ring, a scalar annihilating a nonzero element cannot
be a unit, because multiplication by a unit is injective.
A nilpotent scalar cannot be a unit either, since an invertible
nilpotent would force $1=0$.
Such scalars therefore give nonformality for $r\geq2$.
The criterion permits nilpotents elsewhere in the coefficient ring,
as the positive example over $\mathbb Z/4$ shows.
